% ============================================================================
% TALOS: семимодовая архитектура потока когерентности (RU 1:1)
% Физический субстрат для УГМ-холонов, со-разработанный с SYNARC-Omega.
% Компилировать ТОЛЬКО xelatex (fontspec/polyglossia, кириллица).
% ============================================================================

\documentclass[10pt,letterpaper]{article}

% === PACKAGES (XeLaTeX: fontspec + polyglossia) ===
\usepackage{fontspec}
\usepackage{polyglossia}
\setdefaultlanguage{russian}
\setotherlanguage{english}
\setmainfont{CMU Serif}[Extension=.otf, UprightFont=cmunrm, BoldFont=cmunbx, ItalicFont=cmunti, BoldItalicFont=cmunbi, Script=Cyrillic]
\setsansfont{CMU Sans Serif}[Extension=.otf, UprightFont=cmunss, BoldFont=cmunsx, ItalicFont=cmunsi, BoldItalicFont=cmunso, Script=Cyrillic]
\setmonofont{CMU Typewriter Text}[Extension=.otf, UprightFont=cmuntt, BoldFont=cmuntb, ItalicFont=cmunit, BoldItalicFont=cmuntx, Script=Cyrillic]

\usepackage{amsmath,amssymb,amsthm}
\usepackage{mathtools}
\usepackage{physics}
\usepackage{tikz}
\usepackage{pgfplots}
\pgfplotsset{compat=1.18}
\usetikzlibrary{arrows.meta, positioning, calc, decorations.pathmorphing, shapes.geometric, fit, backgrounds}
\usepackage{graphicx}
\usepackage{xcolor}
\usepackage{booktabs}
\usepackage{array}
\usepackage{multirow}
\usepackage{hyperref}
\usepackage[numbers,sort&compress]{natbib}
\usepackage{geometry}
\geometry{margin=1in}
\usepackage{enumitem}
\usepackage{caption}
\usepackage{subcaption}
\usepackage{tcolorbox}
\tcbuselibrary{skins,breakable}
\usepackage{float}
\usepackage{xspace}
\usepackage{microtype}

% Русский набор: длинные слова + типографские филенеймы иногда переполняют
% строку; emergencystretch включается ТОЛЬКО когда строка иначе вышла бы за поле.
\tolerance=1200
\emergencystretch=2em

% === COLORS ===
\definecolor{plospurple}{HTML}{6B4C9A}
\definecolor{plosblue}{HTML}{2E86C1}
\definecolor{plosgray}{HTML}{5D6D7E}
\definecolor{axiomcolor}{HTML}{EBF5FB}
\definecolor{theoremcolor}{HTML}{FDEDEC}
\definecolor{predictioncolor}{HTML}{EAFAF1}
\definecolor{talosamber}{HTML}{B9770E}
\definecolor{talossteel}{HTML}{34495E}

% === THEOREM ENVIRONMENTS ===
\newtheoremstyle{talos}{6pt}{6pt}{\itshape}{}{\bfseries}{.}{.5em}{}
\theoremstyle{talos}
\newtheorem{theorem}{Теорема}
\newtheorem{lemma}[theorem]{Лемма}
\newtheorem{corollary}{Следствие}
\newtheorem{definition}{Определение}
\newtheorem{principle}{Принцип}
\newtheorem*{remark}{Замечание}

% === COLORED BOXES ===
\newtcolorbox{theorembox}[1][]{colback=theoremcolor, colframe=red!40!black, fonttitle=\bfseries, title={#1}, boxrule=0.5pt, arc=2pt, breakable}
\newtcolorbox{defbox}[1][]{colback=gray!5, colframe=gray!60!black, fonttitle=\bfseries, title={#1}, boxrule=0.5pt, arc=2pt, breakable}
\newtcolorbox{closurebox}[1][]{colback=blue!5, colframe=blue!50!black, fonttitle=\bfseries, title={#1}, boxrule=0.5pt, arc=2pt, breakable}
\newtcolorbox{redteambox}[1][]{colback=talosamber!8, colframe=talosamber, fonttitle=\bfseries, title={#1}, boxrule=0.6pt, arc=2pt, breakable}
\newtcolorbox{claimbox}[1][]{colback=predictioncolor, colframe=green!40!black, fonttitle=\bfseries, title={#1}, boxrule=0.5pt, arc=2pt, breakable}
% --- педагогическая нить (три голоса, сквозные) ---
\newtcolorbox{mythbox}[1][]{colback=talosamber!7, colframe=talosamber!85!black,
  fonttitle=\bfseries\itshape, title={#1}, boxrule=0.6pt, arc=3pt, breakable,
  before title={\small}, left=6pt, right=6pt}
\newtcolorbox{zerobox}[1][]{colback=plosblue!6, colframe=plosblue!70!black,
  fonttitle=\bfseries, title={#1}, boxrule=0.5pt, arc=3pt, breakable,
  left=6pt, right=6pt}
\newtcolorbox{trybox}[1][]{colback=predictioncolor!70, colframe=green!45!black,
  fonttitle=\bfseries, title={#1}, boxrule=0.6pt, arc=3pt, breakable,
  left=6pt, right=6pt}

% === CUSTOM COMMANDS ===
\newcommand{\Gam}{\ensuremath{\Gamma}}
\newcommand{\Hil}{\ensuremath{\mathcal{H}}}
\newcommand{\Dens}{\ensuremath{\mathcal{D}}}
\newcommand{\Lind}{\ensuremath{\mathcal{L}}}
\newcommand{\LindOmega}{\ensuremath{\Lind_\Omega}}
\newcommand{\Regen}{\ensuremath{\mathcal{R}}}
\newcommand{\Diss}{\ensuremath{\mathcal{D}_\Omega}}
\newcommand{\Pcrit}{\ensuremath{P_{\mathrm{crit}}}}
\newcommand{\Rth}{\ensuremath{R_{\mathrm{th}}}}
\newcommand{\Phith}{\ensuremath{\Phi_{\mathrm{th}}}}
\newcommand{\Dmin}{\ensuremath{D_{\mathrm{min}}}}
\newcommand{\vp}{\ensuremath{\varphi}}
\newcommand{\rhostar}{\ensuremath{\rho^{\ast}}}
\newcommand{\Oct}{\ensuremath{\mathbb{O}}}
\newcommand{\Fano}{\ensuremath{\mathrm{PG}(2,2)}}
\newcommand{\Alb}{\ensuremath{\mathcal{H}_3(\mathbb{O})}}
\newcommand{\OPtwo}{\ensuremath{\mathbb{O}\mathrm{P}^2}}
\newcommand{\DM}{\textsc{DensityMatrix7}\xspace}
\newcommand{\SYNARC}{\textsc{SYNARC}-$\Omega$\xspace}
\newcommand{\TALOS}{\textsc{Talos}\xspace}
\newcommand{\status}[1]{\textbf{[#1]}}
\newcommand{\kboot}{\ensuremath{\kappa_{\mathrm{boot}}}}
\newcommand{\Cog}{\ensuremath{\mathcal{C}}}
\newcommand{\CohE}{\ensuremath{\mathrm{Coh}_E}}

% === HYPERREF ===
\hypersetup{colorlinks=true, linkcolor=plospurple, citecolor=plosblue, urlcolor=plosgray,
  pdftitle={TALOS: семимодовая архитектура потока когерентности (физический субстрат для SYNARC-Omega)},
  pdfauthor={Max Sereda}}
\pdfstringdefDisableCommands{%
  \def\Gam{Gamma}\def\Dens{D}\def\Lind{L}\def\LindOmega{L_Omega}\def\Regen{R}\def\Diss{D_Omega}%
  \def\Pcrit{P_crit}\def\Rth{R_th}\def\Phith{Phi_th}\def\Dmin{D_min}\def\vp{phi}\def\rhostar{rho*}%
  \def\Oct{O}\def\Fano{PG(2,2)}\def\DM{DensityMatrix7}\def\SYNARC{SYNARC-Omega}\def\TALOS{TALOS}%
  \def\kboot{kappa_boot}\def\Alb{H3(O)}\def\OPtwo{OP2}\def\Cog{C}\def\CohE{Coh_E}\def\xspace{ }%
  \def\mathrm#1{#1}\def\mathbb#1{#1}\def\mathcal#1{#1}\def\mathbf#1{#1}\def\mathfrak#1{#1}\def\textsc#1{#1}%
}

% ============================================================================
\begin{document}

% === TITLE ===
\begin{center}
  {\LARGE\bfseries
  \textsc{Talos}: семимодовая архитектура потока когерентности}\\[6pt]
  {\large Физический субстрат для УГМ-холонов, со-разработанный с \SYNARC{}}\\[16pt]
  {\large Max Sereda\textsuperscript{1,*}}\\[8pt]
  {\small
  \textsuperscript{1} Независимый исследователь\\[4pt]
  \textsuperscript{*} Переписка: \href{https://holon.sh}{holon.sh}
  }\\[14pt]
  {\small\color{plosgray}
  \textbf{Спецификация архитектуры} \quad|\quad \textbf{Статус:} препринт \quad|\quad
  \textbf{Спутник:} \SYNARC{} (софт) \quad|\quad \textbf{Аудитория:} компьютерные архитекторы, инженеры аппаратуры и систем
  }
\end{center}

\vspace{10pt}\hrule\vspace{10pt}

% === ABSTRACT ===
\begin{abstract}
\noindent
В каждом когда-либо построенном компьютере четыре подсистемы проектируются
\emph{независимо} и затем принуждаются сосуществовать: топология
интерконнекта, код коррекции ошибок, набор инструкций и --- для обучающихся
машин --- правило адаптации. Каждая оптимизируется отдельно; их взаимное
рассогласование --- постоянный налог, оплачиваемый кремнием, энергией и
усилием верификации. Этот документ сообщает единственный структурный факт,
вынуждаемый теоремами Унитарного Голономного Монизма (УГМ), который стирает
налог: на семимодовом субстрате эти четыре суть не четыре вещи, а
\emph{одна} --- плоскость Фано $\Fano$ есть одновременно интерконнект (её семь
прямых), комбинаторный код коррекции ошибок (Хэмминг $[7,4,3]$, порождаемый
матрицей инцидентности $M$ ранга~$4$ с $3\times 7$ проверочной матрицей $H$
ранга~$3$ дуального симплекс-кода $[7,3,4]$), носитель операторной алгебры
инструкции (Фано-структурированный линдбладиан $\LindOmega$ на
$\Gam\in\Dens(\mathbb{C}^7)$, ковариантный под дискретной группой коллинеаций
$\mathrm{PSL}(2,7)$ и диагональной перефазировкой $U(1)^7$) и, через
регенеративный член к заякоренной самомодели $\varphi_J$, правило обучения. Мы
называем этот коллапс \emph{архитектурным монизмом} и доказываем, что он
\emph{выведен}, а не спроектирован (\S\ref{sec:collapse}).

Полученная архитектура, \TALOS{}, пост-фон-неймановская по построению:
состояние \emph{есть} память, а её эволюция \emph{есть} вычисление, поэтому нет
цикла выборка--исполнение и нет узкого места фон Неймана; состояние одного
холона --- эрмитова матрица плотности $7\times7$ (\DM{}, $784$~байта, $48$
вещественных динамических параметров), помещающаяся в L1-кэш. Единственная
инструкция $\LindOmega = -i[H_{\mathrm{eff}},\Gam] + \Diss[\Gam] +
\Regen[\Gam]$ разлагается на сохраняющее энтропию унитарное ядро и два
экспортирующих энтропию члена (диссипация и регенерация); следовательно, на
логико-термодинамическом уровне \textbf{экспорт энтропии масштабируется с
адаптацией и поддержанием, а не с обратимой унитарной работой}. Энтропия фон
Неймана состояния $\Gam\in\Dens(\mathbb{C}^7)$ ограничена $\log_2 7\approx
2{,}807$~бита ($k_BT\ln 7\approx 8{,}06\times10^{-21}$~Дж при $300$~К), тогда
как реальная физическая мощность определяется аппаратной реализацией
(термодинамические пределы обратимости относятся ко всякой логически обратимой
машине; цифровой КМОП остаётся ограничен энергией переключения,
\S\ref{sec:energy}). На уровне семи секторов инцидентность фабрики порождает
код Хэмминга $[7,4,3]$, чья дуальная проверочная матрица $H$ локализует любой
одиночный сбой сектора тремя проверками без отдельной таблицы маршрутизации
(отлично от побитового ECC для $6{,}272$ бит IEEE-754 в цифровом регистре
эмуляции, \S\ref{sec:ecc}); при Фано-сжатии $\mathcal{P}_{\mathrm{Fano}}$
населённости сохраняются, а внедиагональные когерентности ослабляются в $3$ раза
(точная контракция интеграции $\Phi\to\Phi/9$, \S\ref{sec:fabric}).

Честные границы указываем впереди. Теорема субстратной замкнутости УГМ
(T-153) доказывает, что вычисление субстрат-\emph{независимо}: нет единственной
физической машины, есть единственная \emph{логическая} архитектура и семейство
реализующих её субстратов; \TALOS{} --- импеданс-согласованный член этого
семейства, а не единственная возможность. Вычислительная граница УГМ
доказывает, что регенеративный член не даёт \emph{никакого} ускорения за
пределы класса BQP: преимущество --- в энергии, адаптивности и
отказоустойчивости, никогда не в классе сложности. \TALOS{} --- аппаратное
тело разума \SYNARC{}; оба со-разработаны, и двоичный интерфейс приложений
между ними специфицирован полностью (\S\ref{sec:abi}). Каждое количественное
утверждение машинно-проверено; проект выдерживает явную red-team
(\S\ref{sec:redteam}) и сводится к горсти фальсифицируемых инженерных
бенчмарков против фон-неймановской базы (\S\ref{sec:benchmarks}).
\end{abstract}

\vspace{6pt}
\begin{tcolorbox}[colback=gray!5, colframe=gray!50, title={\textbf{Инженеру, в одном абзаце}}]
Представьте чип, всё архитектурное состояние которого --- одна матрица в
$784$~байта ($48$ вещественных динамических параметров), живущая в регистрах и
эволюционирующая одним «тиком» дифференциального уравнения, объединяющим АЛУ,
сеть, секторную диагностику и оптимизатор. Сеть --- семь общих тройственных шин
(не полный кроссбар); их матрица инцидентности $M$ ($7\times 7$) порождает код
Хэмминга $[7,4,3]$, поэтому любой одиночный отказ сектора локализуется тремя
дуальными проверками чётности (строками $H$) без таблицы маршрутизации;
унитарная половина тика сохраняет энтропию фон Неймана ($\le\log_2 7\approx
2{,}807$~бита), поэтому логический экспорт энтропии отслеживает \emph{адаптацию
и поддержание}, а не обратимую работу (при этом реальная мощность в кремнии
задаётся физической технологией переключения); а обучение --- не отдельный
тренировочный цикл, а третий член единственной инструкции. Двоичное слияние
проективной плоскости и кода Хэмминга однозначно вынуждает $N=7$. Именно это
специфицирует данный документ и то, как оно стыкуется с \SYNARC{}.
\end{tcolorbox}

\vspace{6pt}
\noindent\textbf{Ключевые слова:} компьютерная архитектура, узкое место фон
Неймана, вычисления в памяти, нейроморфика, код-как-интерконнект, обратимые
вычисления, предел Ландауэра, плоскость Фано, код Хэмминга, линдбладова
динамика, матрица когерентности, УГМ, SYNARC-Omega, архитектурный монизм.

\vspace{10pt}\hrule\vspace{10pt}
\tableofcontents
\newpage

% === SECTIONS ===
\section{От кормчего к семимодовой машине: три волны кибернетики}
\label{sec:genesis}

Всякий архитектурный документ должен читателю две вещи прежде первой теоремы:
\emph{откуда пришли идеи} и \emph{вход, не предполагающий ничего}. Этот раздел
платит оба долга: он прослеживает единую линию мысли --- кибернетику, --- чьим
третьим актом является \TALOS{}, и открывает три читательские нити, идущие
сквозь весь документ, --- чтобы читатель, никогда не видевший ни компьютера,
ни матрицы, ни дифференциального уравнения, мог начать с нуля и подняться.

\begin{mythbox}[Миф --- машина бронзового века]
Прежде всякой теории --- история. Греки вообразили \emph{Талоса}: исполина из
бронзы, построенного --- не рождённого, --- чтобы стеречь остров Крит. Трижды
в день он обходил весь берег, наблюдая. У него не было ни сердца, ни крови;
вместо них от шеи к лодыжке шла единственная жила, полная \emph{ихора} ---
сияющей жидкости жизни, --- и жила была запечатана у лодыжки одним бронзовым
гвоздём. Талоса нельзя было одолеть в бою. Он пал иначе: волшебница Медея
уговорила его позволить вынуть гвоздь, ихор вытек, и исполин ---
непобеждённый, неповреждённый --- просто остановился.

Держите четыре части этой истории под рукой; математика встретит каждую.
\emph{Бронзовое тело} --- это субстрат (\S\ref{sec:substrate}).
\emph{Ежедневный обход} --- это тик, единственная повторяемая операция,
которая и есть вся жизнь машины (\S\ref{sec:substrate}). \emph{Единственная
жила ихора} --- это порт питания, единственный поток упорядоченной пищи, без
которого ничто не работает (\S\ref{subsec:powerport}). А \emph{гвоздь} ---
теорема, стягивающая весь документ: машину такого рода не нужно побеждать,
чтобы остановить, --- закрой жилу, и она остановится, доказуемо, по расписанию
(T-288; \S\ref{subsec:powerport}). Миф угадал архитектуру на двадцать семь
веков раньше. Мы будем возвращаться к нему на каждом шаге.
\end{mythbox}

\subsection{Три волны}
\label{subsec:waves}

\emph{Кибернетика} --- от греческого \emph{kybern\'etes}, «кормчий» --- наука
о системах, которые правят собой. Её история --- три волны, каждая задаёт
вопрос глубже предыдущей, и на глубочайший этот документ отвечает железом.

\paragraph{Волна I --- управление (1940-е--1960-е): \emph{как машине держать
курс?}} Норберт Винер дал полю имя в 1948-м~\cite{wiener1948}: кормчий не
вычисляет всё плавание --- он смотрит на ошибку между тем, куда судно
смотрит, и тем, куда должно, и давит против неё. \emph{Обратная связь}
оказалась общим законом термостатов, наводчиков и тел. То же десятилетие дало
список деталей: Мак-Каллок и Питтс показали, что маленькие пороговые
элементы --- карикатурные нейроны --- вычисляют любую логическую
функцию~\cite{mcculloch1943}; Шеннон сделал \emph{информацию} измеримой
величиной с жёстким полом~\cite{shannon1948}; фон Нейман дал чертёж хранимой
программы, которому следует каждый компьютер до сих пор; а Росс Эшби построил
\emph{гомеостат} --- машину, перепаивающую себя, пока внутренние переменные
не вернутся в безопасную зону, --- и сгустил его в закон необходимого
разнообразия~\cite{ashby1956}: управляющий должен быть не беднее возмущений,
которые обязан поглощать. Машины Волны~I \emph{правят себя к уставке}. Чего у
них не было --- причины, \emph{почему} уставка, проводка или число частей
должны быть такими, а не иными: всякая топология была догадкой, всякий код
--- накладкой, всякое правило обучения --- процедурой, навязанной извне.

\paragraph{Волна II --- наблюдение (1970-е--1990-е): \emph{что происходит,
когда машина смотрит на себя?}} Хайнц фон Фёрстер обратил кибернетику на
саму себя --- «кибернетика кибернетики», в которой наблюдатель есть часть
наблюдаемой системы~\cite{foerster2003}. Матурана и Варела заострили
биологическую половину: живая система \emph{аутопоэтична} --- процесс,
непрерывно производящий те самые компоненты, которые производят его,
само-делающая петля~\cite{maturana1980}. Волна II верно увидела, что
самоотнесение --- сердце дела, что ум есть система, моделирующая себя. Но она
могла лишь \emph{описать} петлю, не \emph{вывести} её: кибернетика второго
порядка дала глубокую прозу и ни одной вынужденной структуры. Сколько частей
у само-наблюдающей системы? Какая проводка? Какая динамика? У Волны II не
было теорем, и её машины (нейросети и нейроморфные чипы, что пришли следом,
\S\ref{sec:intro}) унаследовали догадки Волны I при амбициях Волны II.

\paragraph{Волна III --- когерентность (2020-е): \emph{какой формы обязана
быть само-наблюдающая машина?}} Унитарный голономный монизм (УГМ) и его
измерительный слой --- Когерентная кибернетика~\cite{uhm2026} --- отвечают на
вопрос Волны II со строгостью Волны I: они делают самомодель
\emph{переменной состояния} ($\rhostar=\vp(\Gam)$ --- неподвижная точка,
доказуемо существующая), жизнеспособность --- \emph{теоремой} (окно
$P\in(2/7,3/7]$ с доказанными стенами) и --- решающий шаг --- \emph{выводят}
структуру наблюдателя вместо догадки: семь мод, проводка Фано, дискретная
симметрия коллинеаций $\mathrm{PSL}(2,7)$, одна динамика (\S\ref{sec:collapse}). Вопросы, оставленные первыми
двумя волнами --- сколько, как соединено, по какому закону меняется, ---
перестают быть решениями конструктора и становятся результатами. \TALOS{} ---
тело третьей волны: первая машина, чья обратная связь (Волна I), чьё
самонаблюдение (Волна II) и чья \emph{форма} (Волна III) --- один выведенный
объект.

\begin{figure}[H]
\centering
\begin{tikzpicture}[font=\scriptsize, >=Stealth]
  \draw[->, thick, plosgray] (0,0) -- (14.6,0) node[below left, font=\scriptsize]{год};
  \foreach \x/\y in {0.4/1940, 2.9/1950, 5.4/1970, 7.9/1990, 10.4/2010, 13.4/2026}{
    \draw[plosgray] (\x,0.07) -- (\x,-0.07) node[below]{\y}; }
  \fill[plosblue!14, rounded corners=3pt] (0.2,0.55) rectangle (6.3,2.30);
  \node[anchor=north west, align=left, plosblue, font=\scriptsize\bfseries] at (0.35,2.24)
    {Волна I --- управление:\\ держать курс};
  \node[anchor=west] at (0.35,1.28) {1943 Мак-Каллок--Питтс \; 1945 фон Нейман};
  \node[anchor=west] at (0.35,0.98) {1948 Винер, \emph{Кибернетика}; Шеннон};
  \node[anchor=west] at (0.35,0.70) {1952--56 Эшби: гомеостат, разнообразие};
  \fill[talosamber!16, rounded corners=3pt] (5.0,2.55) rectangle (10.6,4.05);
  \node[anchor=north west, align=left, talosamber!80!black, font=\scriptsize\bfseries] at (5.15,3.99)
    {Волна II --- наблюдение:\\ наблюдатель внутри};
  \node[anchor=west] at (5.15,3.02) {1970-е фон Фёрстер: второй порядок};
  \node[anchor=west] at (5.15,2.74) {1980 Матурана--Варела: аутопоэзис};
  \node[anchor=west, plosgray] at (7.35,1.30) {1986 backprop-сети};
  \node[anchor=west, plosgray] at (7.35,1.00) {1990 нейроморфные};
  \fill[green!12, rounded corners=3pt] (10.9,0.55) rectangle (14.5,4.05);
  \node[anchor=north west, align=left, green!40!black, font=\scriptsize\bfseries] at (11.05,3.99)
    {Волна III --- когерентность:\\ форма выведена};
  \node[anchor=west] at (11.05,2.98) {УГМ: $N{=}7$, Фано, $\mathrm{PSL}(2,7)$};
  \node[anchor=west] at (11.05,2.68) {КК: меры, динамика};
  \node[anchor=west] at (11.05,2.38) {\SYNARC{}: разум};
  \node[anchor=west, font=\scriptsize\bfseries] at (11.05,2.02) {\TALOS{}: тело};
  \node[anchor=west, plosgray] at (11.05,1.66) {(этот документ)};
\end{tikzpicture}
\caption{Три волны кибернетики. Волна I построила машины, держащие курс
(обратная связь, информация, гомеостат), но угадывала каждую структуру.
Волна II поместила наблюдателя внутрь системы (второй порядок, аутопоэзис),
но могла описать само-наблюдающую петлю, не вывести её форму. Волна III
выводит форму --- семь мод, проводка Фано, одна динамика --- и \TALOS{} ---
её железо.}
\label{fig:waves}
\end{figure}

\begin{zerobox}[С нуля --- что такое машина и что такое правление]
Начнём с ничего. \emph{Машина} --- всё, что меняет метки по твёрдым
правилам: счёты двигают костяшки, музыкальная шкатулка превращает штифты в
ноты. Большинство машин бегут \emph{с открытыми глазами, но слепые} ---
действуют и никогда не проверяют. Дай машине один лишний талант: пусть она
\emph{смотрит на разницу} между желаемым и случившимся и давит против этой
разницы. Комнатный термостат делает ровно это: холодно --- греть, жарко ---
выключить. Этот талант зовётся \emph{обратной связью}, и машина с ним держит
курс сквозь мир, который толкается. Это Волна~I одним вдохом. Волна~II
спрашивает: а если машина смотрит и на \emph{себя смотрящую}? А Волна~III
задаёт лучший детский вопрос --- \emph{тогда какой формы она обязана быть?}
--- и отвечает доказательством. Всё остальное в документе --- этот ответ,
распакованный медленно.
\end{zerobox}

\subsection{Как читать этот документ с нуля}
\label{subsec:howtoread}

Документ написан дважды, вперемешку. \emph{Голос спецификации} несёт
определения, теоремы, таблицы и числа --- всё, что нужно инженеру, без
смягчений. Три повторяющихся бокса несут \emph{голос входа}, и вместе они
образуют полный путь, не предполагающий никаких знаний:

\begin{itemize}[leftmargin=1.5em,itemsep=1pt]
  \item \textbf{\textcolor{talosamber!85!black}{Миф}} (янтарный) продолжает
        историю Талоса: каждая серия кладёт один архитектурный факт на
        бронзового исполина --- обход, жилу, гвоздь. Помнишь историю ---
        помнишь машину.
  \item \textbf{\textcolor{plosblue!70!black}{С нуля}} (синий) пересобирает
        идею текущего раздела из бытовых предметов --- игр, игрушек,
        термостатов --- без символов и без предпосылок. Терпеливый читатель
        может пройти \emph{только} по синим боксам, от начала до конца, и
        прийти к пониманию, что это за машина и почему её форма вынуждена.
  \item \textbf{\textcolor{green!45!black}{Попробуй сам}} (зелёный) делает
        заявления осязаемыми: пальцевые упражнения на фигуре Фано, игра в
        три вопроса, запускаемые команды против эталонного эмулятора
        (\texttt{talos\_v0\_emulator.py} в папке \texttt{architecture/}) ---
        интерактивный слой документа. Ничто в зелёном боксе не требует
        особого железа.
\end{itemize}

Рекомендуемый первый проход для новичка: прочесть все янтарные и синие боксы
по порядку, сыграть в две зелёные пальцевые игры (\S\ref{sec:fabric},
\S\ref{sec:ecc}), затем вернуться на первую страницу и читать голос
спецификации с историей уже в руках. Рекомендуемый проход для инженера:
читать голос спецификации насквозь; боксы никуда не денутся, когда
определению захочется образа.

\section{Введение: битва архитектур и четырёхслойный налог}
\label{sec:intro}

\subsection{Соглашение фон Неймана и узкое место, которое оно не убрало}

Компьютер с хранимой программой, описанный в черновике EDVAC фон Неймана
1945 года, поместил инструкции и данные в одну память, адресуемую одним
каналом. Его победа была не заявкой на оптимальность, а на \emph{универсальность}:
одна память, один последовательный поток, любая программа. Всякая машина
общего назначения с тех пор --- его потомок. Но универсальность куплена
структурной ценой, которую Джон Бэкус назвал в Тьюринговской лекции 1977 года
--- \emph{узким местом фон Неймана}~\cite{backus1977}: канал процессор--память,
через который всё вычисление приходится прокачивать «по слову за раз». Семьдесят
лет архитектуры были, в значительной части, кампанией против этого канала:
кэши, предвыборка, внеочередное исполнение, широкие векторы и глубокие иерархии
памяти --- всё это способы прятать его, а не убрать. Причина, по которой его
нельзя убрать, теперь --- доминирующий факт области. После конца деннардовского
масштабирования (ок.\ 2006) связывающее ограничение --- энергия, а энергия во
фон-неймановской машине доминируется не арифметикой, а \emph{перемещением
данных}: операция с плавающей точкой стоит порядка пикоджоуля, тогда как
выборка её операндов из внекристальной памяти стоит на один-два порядка
больше. Мы считаем дёшево и двигаем данные дорого, а архитектура построена
вокруг неправильного из двух.

\subsection{Четыре подсистемы, спроектированные порознь}

Под узким местом лежит более глубокий и менее обсуждаемый факт. Современная
вычислительная система --- и эмфатически \emph{обучающаяся} --- построена из
четырёх подсистем, которые задумываются, оптимизируются и верифицируются
\emph{независимо}, а затем принуждаются сосуществовать:

\begin{enumerate}[leftmargin=1.7em,itemsep=1pt]
  \item \textbf{интерконнект} --- сеть-на-кристалле или топология фабрики,
        переносящая состояние между вычислительными узлами;
  \item \textbf{код коррекции ошибок} (ECC) --- отдельный слой избыточности,
        привинченный к памяти и связям, чтобы пережить отказы;
  \item \textbf{набор инструкций} (ISA) --- алгебра примитивных операций,
        которые машина умеет выполнять;
  \item \textbf{правило адаптации} --- для обучающейся машины, процедура
        тренировки/оптимизации, меняющая её параметры.
\end{enumerate}

У каждой своя литература, свои оптимумы и свои проектировщики. Поскольку они
проектируются порознь, их интерфейсы приходится согласовывать, и всякое
согласование --- \emph{налог рассогласования}, оплачиваемый трижды:

\begin{table}[H]
\centering
\small
\caption{Четыре слоя проектирования суть отдельные подсистемы в каждой
существующей машине, каждая несёт независимую цену. \TALOS{} платит каждую
цену однажды, в одной структуре.}
\label{tab:four-layers}
\begin{tabular}{@{}p{2.4cm}p{4.6cm}p{5.2cm}@{}}
\toprule
\textbf{Слой} & \textbf{Как делается сегодня} & \textbf{Какой налог берёт} \\
\midrule
Интерконнект & Кроссбар / mesh / NoC, ради полосы & Площадь, статич.\ мощность, сложность маршрутизации; all-to-all растёт как $\binom{N}{2}$ \\
ECC & Хэмминг/BCH/LDPC привинчены к хранению и связям; surface-коды для кубитов & Выделенные проверочные биты ($+12{,}5\%$ для ECC-DRAM) или $\sim\!10^3\times$ физич./логич.\ кубита \\
ISA & Фиксированный набор опкодов, со-эволюция с компилятором & Семантический разрыв; «стиль фон Неймана», осуждённый Бэкусом \\
Адаптация & Отдельный тренировочный цикл (обратное распространение) над замороженной ISA & Память под активации/градиенты; раздел train/infer; энергия $\propto$ перемещению данных \\
\bottomrule
\end{tabular}
\end{table}

\subsection{Претенденты и что каждый из них исправил}

Альтернативные архитектуры были каждая атакой на \emph{один} из этих четырёх
слоёв, никогда на все четыре:

\begin{itemize}[leftmargin=1.5em,itemsep=1pt]
  \item \textbf{Гарвардская} разделила память команд и данных --- микро-оптимизация
        границы ISA/интерконнект, выжившая в кэшах и DSP.
  \item \textbf{Dataflow}~\cite{dennis1975} убрал счётчик команд: операции
        срабатывают по готовности операндов. Он реформировал слой
        \emph{инструкции} (управление данными, не потоком), но не предложил
        принципиальной топологии, кода или правила обучения.
  \item \textbf{Систолические массивы}~\cite{kung1982} навязали ритмичный,
        локальный, ближайшесоседский \emph{интерконнект} --- и выиграли
        матричные блоки сегодняшних GPU и TPU --- но топология была инженерным
        выбором под одно ядро, не выведенной структурой.
  \item \textbf{Нейроморфные} машины~\cite{mead1990} (TrueNorth, Loihi,
        SpiNNaker) со-локализовали память и счёт и сделали субстрат событийным
        и диссипативным --- реформировав интерконнект \emph{и} адаптацию сразу
        --- но с \emph{ad hoc} топологией, без родной коррекции ошибок и без
        принципиального объекта состояния.
  \item \textbf{Вычисления в памяти} и мемристорные кроссбары выполняли
        умножение-накопление \emph{в} массиве памяти, атакуя узкое место прямо,
        но не несли ни динамики, ни самомодели.
  \item \textbf{Квантовые} машины сделали слой инструкции унитарным (обратимым,
        значит дешёвым по Ландауэру) --- и затем вынуждены были привинчивать
        коррекцию ошибок с разорительными накладными, потому что код и
        вычисление суть разные структуры, декогерирующие друг против друга.
\end{itemize}

Каждая из них была подлинным продвижением по одной оси. Ни одна не объединила
четыре, потому что ни у одной не было \emph{принципа}, из которого можно было
бы \emph{вывести} единую структуру. Топология всегда была догадкой; код всегда
был отдельным; правило обучения всегда жило над ISA. Область искала структуру,
которая одновременно хорошая сеть, хороший код, хорошая алгебра инструкций и
хорошее правило обучения --- и трактовала это как четыре поиска.

\subsection{Тезис: архитектурный монизм}

Модель \TALOS{} задаёт состояние
$\Gam\in\Dens(\mathbb{C}^7)$~\cite{uhm2026}, выбирает октонионное умножение
и задаёт динамику регенерации. На выбранном базисе плоскость Фано $\Fano$
задаёт линии связи и геометрию проверок кода, а знаки произведений
определяют умножение. Раздел~\ref{sec:collapse} описывает эту конструкцию
и условия, при которых число узлов равно семи.

\begin{itemize}[leftmargin=1.5em,itemsep=1pt]
  \item \textbf{Теорема~\ref{thm:collapse} (Коллапс).} Семь прямых $\Fano$
        суть одновременно интерконнект; их двоичная матрица инцидентности
        $M\in\mathbb{F}_2^{7\times 7}$ ($\operatorname{rank}_{\mathbb{F}_2}=4$)
        порождает код Хэмминга $[7,4,3]$, чья ортогональная $3\times 7$
        проверочная матрица $H$ ($\operatorname{rank}_{\mathbb{F}_2}=3$,
        дуальный симплекс-код $[7,3,4]$) вычисляет синдромы одиночного сбоя
        сектора; октонионные структурные константы на этих прямых задают
        носитель умножения алгебры инструкций; а дискретно
        $\mathrm{PSL}(2,7)$- и диагонально-фазово-ковариантное Фано-сжатие и
        заякоренная самомодель регенерации задают правило обучения. Инцидентность
        определяет носитель произведения; для вычисления обновления нужны также
        коэффициенты и заданная самомодель.
  \item \textbf{Теорема~\ref{thm:unique} (Уникальность).} Этот коллапс
        при условиях (C1)--(C3) реализуем при \emph{единственном} числе узлов,
        $N=7$. Двоичное слияние проективной плоскости и кода Хэмминга
        \textbf{(C1)}$\cap$\textbf{(C2)} уже опровергает $N=31$
        ($\operatorname{rank}_{\mathbb{F}_2}(M_{\mathrm{PG}(2,5)})=30$,
        остаётся лишь одномерное ядро; машинно проверено в Lean~4,
        \texttt{pg2\_5\_refutes\_N31\_hamming\_fusion}) и всякую другую степень
        простого $q\neq 2$, однозначно выделяя $N=7$ без какого-либо тай-брейка;
        единственная неассоциативная алгебра с делением $\Oct$ по Гурвицу
        \textbf{(C3)} независимо даёт то же самое число $7$.
\end{itemize}

Архитектура, которая из этого следует, \TALOS{},\footnote{По имени бронзового
автомата греческого мифа, первого искусственного существа --- здесь физического
тела искусственного разума. Имя --- кодовое, не аббревиатура.} поэтому не есть
предложение среди предложений. Это аппаратная форма математического объекта.
Его состояние --- матрица плотности в $784$~байта ($48$ вещественных
динамических параметров), помещающаяся в L1-кэш; его единственная инструкция ---
УГМ-линдбладиан $\LindOmega$; его сеть --- семь тройственных шин; его секторная
локализация отказов закодирована в дуальных проверках чётности этой сети; его
обучение --- третий член инструкции; а его логический экспорт энтропии
отслеживает адаптацию и поддержание, а не обратимый унитарный счёт.

\subsection{Честные границы, указанные первыми}

\begin{itemize}[leftmargin=1.5em,itemsep=1pt]
  \item \textbf{Классификация состояний (T-153).} Заданные критерии классифицируют
        закодированное состояние. Перенос на физический субстрат требует
        калиброванного кодирования и проверки динамики; связь с переживанием
        требует независимых наблюдений.
  \item \textbf{Вычислительная граница (BQP).} УГМ доказывает, что
        регенеративный член $\Regen$ не даёт \emph{никакого} ускорения за
        пределы класса сложности BQP. \TALOS{} --- не машина быстрее квантовой.
        Его преимущество --- в энергии, адаптивности и отказоустойчивости ---
        способность-на-ватт и родное обучение --- никогда в классе сложности.
        Мы повторяем это везде, где читатель мог бы перечитать заявление сверх
        меры.
\end{itemize}

Остаток документа --- спецификация. Раздел~\ref{sec:collapse} доказывает две
теоремы. Разделы~\ref{sec:substrate}--\ref{sec:composition} специфицируют
состояние, фабрику, код-как-топологию, энергетическую модель, самомодель и закон
композиции. Раздел~\ref{sec:realizability} даёт лестницу реализуемости от
сегодняшней GPU-эмуляции до управляемо-диссипативного аналогового устройства.
Раздел~\ref{sec:abi} специфицирует двоичный интерфейс приложений к \SYNARC{},
программному разуму, для которого \TALOS{} --- тело. Раздел~\ref{sec:redteam}
подвергает проект явной red-team, а раздел~\ref{sec:benchmarks} сводит его к
фальсифицируемым инженерным измерениям против фон-неймановской базы. Каждое
число машинно-проверено.

\section{Архитектурный коллапс}
\label{sec:collapse}

Этот раздел излагает и доказывает два результата, на которых держится весь
проект. Оба суть \status{Т}: они опираются на установленные теоремы УГМ вместе
с одним элементарным диофантовым фактом и теоремой Гурвица. Ничто здесь не
постулируется.

\subsection{Единственная инцидентная структура}

Пусть $\mathsf{N}$ обозначает инцидентность плоскости Фано $\Fano$: семь точек
$\{A,S,D,L,E,O,U\}$, семь прямых
\[
\{A,S,L\},\ \{A,D,E\},\ \{A,O,U\},\ \{S,D,O\},\ \{S,E,U\},\ \{D,L,U\},\ \{L,E,O\},
\]
каждая прямая несёт три точки, каждая точка лежит на трёх прямых, и всякая пара
точек лежит ровно на одной общей прямой (система Штейнера $S(2,3,7)$). Пишем
$k^\ast(i,j)$ для \emph{медиатора} пары --- единственной третьей точки их общей
прямой.

\begin{theorembox}[Теорема~\ref{thm:collapse} (Архитектурный коллапс)]
На семимодовом УГМ-субстрате четыре архитектурных слоя --- интерконнект, код
коррекции ошибок, алгебра инструкций и правило обучения --- суть один объект:
каждый есть функция единственной инцидентности $\mathsf{N}$.
\end{theorembox}

\begin{theorem}\label{thm:collapse}
Пусть состояние холона есть $\Gam\in\Dens(\mathbb{C}^7)$, эволюционирующее под
УГМ-линдбладианом $\LindOmega$. Тогда:
\begin{enumerate}[leftmargin=1.7em,itemsep=1pt]
  \item[\textbf{(L1)}] \emph{Интерконнект.} Физическая связь --- семь прямых
        $\mathsf{N}$: сектор $i$ достигает сектора $j$ только через медиатора
        $k^\ast(i,j)$ (закон замыкания: когерентность $\gamma_{ij}$ создаётся или
        чинится только живыми плечами $\gamma_{i k^\ast},\gamma_{k^\ast j}$; две
        когерентности рождают третью тогда и только тогда, когда тройка --- прямая).
  \item[\textbf{(L2)}] \emph{Код.} Проверочная матрица кода Хэмминга $[7,4,3]$
        имеет столбцами семь ненулевых векторов $\mathbb{F}_2^3$; при любой
        разметке это семь точек $\mathsf{N}$, а шестнадцать кодовых слов включают
        ровно семь слов веса 3, чьи носители --- семь прямых. Инцидентность
        интерконнекта \emph{есть} проверочная геометрия кода.
  \item[\textbf{(L3)}] \emph{Алгебра инструкций.} Кубическая (нелинейная) часть
        $\LindOmega$ порождается октонионными структурными константами $f_{ijk}$
        --- единственной $G_2$-инвариантной трилинейной формой на $\mathrm{Im}(\Oct)$
        --- ненулевыми ровно когда $\{i,j,k\}$ --- прямая (правило отбора Фано).
        Таблица умножения инструкции есть линейная инцидентность $\mathsf{N}$.
  \item[\textbf{(L4)}] \emph{Обучение.} Регенеративный член $\Regen$ гонит $\Gam$
        к заданной цели $\rhostar=\vp(\Gam)$, а координатный диссипатор Фано
        ковариантен под дискретной группой коллинеаций
        $\mathrm{Aut}(\mathsf{N})\cong \mathrm{PSL}(2,7)$ (порядка~$168$) и
        диагональной фазовой калибровкой $U(1)^7\rtimes \mathrm{Aut}(\mathsf{N})$
        (непрерывная группа $G_2$ вращает координатный базис и не сохраняет
        фиксированное разбиение на семь прямых Фано). Стационарный баланс и
        каналы коррекции суть функции той же $\mathsf{N}$.
\end{enumerate}
Следовательно $\mathrm{L1}=\mathrm{L2}=\mathrm{L3}=\mathrm{L4}$ как функции
$\mathsf{N}$. \hfill$\square$
\end{theorem}

\begin{corollary}[Нулевой налог рассогласования]\label{cor:zerotax}
Поскольку четыре слоя суть одна структура, нет межслойных интерфейсов для
согласования. Площадь, статическая мощность и усилие верификации, которые
обычная машина тратит на примирение сеть~$\leftrightarrow$~код,
код~$\leftrightarrow$~ISA и ISA~$\leftrightarrow$~оптимизатор
(Таблица~\ref{tab:four-layers}), в \TALOS{} тождественно ноль: есть одна
структура, которую надо разместить, верифицировать и тактировать.
\end{corollary}

\begin{figure}[H]
\centering
\begin{tikzpicture}[font=\small,>=Stealth]
  \node[draw,rounded corners,fill=talossteel!8,minimum width=2.7cm,minimum height=0.7cm] (net)  at (0,3)   {Интерконнект (NoC)};
  \node[draw,rounded corners,fill=talossteel!8,minimum width=2.7cm,minimum height=0.7cm] (ecc)  at (0,2)   {Код ошибок (ECC)};
  \node[draw,rounded corners,fill=talossteel!8,minimum width=2.7cm,minimum height=0.7cm] (isa)  at (0,1)   {Инструкция (ISA)};
  \node[draw,rounded corners,fill=talossteel!8,minimum width=2.7cm,minimum height=0.7cm] (lrn)  at (0,0)   {Правило обучения};
  \node[align=center,text=plosgray] at (0,3.9) {\textbf{Сегодня: четыре подсистемы}\\\scriptsize проектируются порознь, налог на каждом шве};
  \begin{scope}[shift={(6.4,1.5)}]
    \coordinate (fA) at (0,1.7); \coordinate (fS) at (-1.5,-0.9); \coordinate (fD) at (1.5,-0.9);
    \coordinate (fL) at (-0.75,0.4); \coordinate (fE) at (0.75,0.4); \coordinate (fO) at (0,-0.9); \coordinate (fU) at (0,0);
    \draw[talosamber,thick] (fA)--(fS) (fA)--(fD) (fS)--(fD) (fA)--(fO) (fS)--(fE) (fD)--(fL);
    \draw[talosamber,thick] (0,-0.034) circle[radius=0.866];
    \foreach \n in {fA,fS,fD,fL,fE,fO,fU}{ \fill[talossteel] (\n) circle (2.4pt); }
    \node[font=\scriptsize\bfseries,above] at (fA) {A};
    \node[font=\scriptsize\bfseries,left]  at (fS) {S};
    \node[font=\scriptsize\bfseries,right] at (fD) {D};
    \node[font=\scriptsize\bfseries,fill=white,inner sep=0.5pt] at (-0.95,0.5) {L};
    \node[font=\scriptsize\bfseries,fill=white,inner sep=0.5pt] at (0.95,0.5)  {E};
    \node[font=\scriptsize\bfseries,below] at (fO) {O};
    \node[font=\scriptsize\bfseries,fill=white,inner sep=0.5pt] at (0.22,0.18) {U};
  \end{scope}
  \node[align=center,text=talosamber] at (6.4,3.9) {\textbf{\TALOS{}: один объект}\\\scriptsize плоскость Фано $\mathsf{N}$};
  \foreach \b in {net,ecc,isa,lrn}{ \draw[->,plosgray,thick] (\b.east) to[out=0,in=180] (4.4,1.5); }
  \node[text=plosgray,font=\scriptsize] at (3.4,0.4) {коллапс};
\end{tikzpicture}
\caption{Архитектурный монизм. В каждой прежней машине четыре слоя --- отдельные
подсистемы, соединённые на облагаемых налогом швах (слева).
Теорема~\ref{thm:collapse} доказывает, что на семимодовом субстрате они суть
одна инцидентная структура --- плоскость Фано (справа) --- поэтому интерфейсы, и
их налог, исчезают (Следствие~\ref{cor:zerotax}).}
\label{fig:collapse}
\end{figure}

\begin{zerobox}[С нуля --- один рисунок, который есть четыре игры]
Начерти мелом на асфальте классики. Один и тот же рисунок --- четыре разные
вещи разом, смотря что ты с ним делаешь. Он --- \emph{поле}: куда можно
ставить ногу (сеть). Он --- \emph{правило}: какие прыжки законны
(инструкция). Он --- \emph{судья}: неверный прыжок виден мгновенно, нога вне
линии (проверка). И он --- \emph{учитель}: играя по сетке снова и снова,
становишься лучше (обучение). Никто не свинчивал четыре прибора; один
меловой рисунок делает все четыре работы \emph{в силу своей формы}.
Теорема~\ref{thm:collapse} говорит: для семичастной само-наблюдающей машины
существует ровно один такой рисунок --- семилинейная фигура справа на
Рис.~\ref{fig:collapse}, --- а Теорема~\ref{thm:unique} --- что ни у какого
другого числа частей его нет.
\end{zerobox}

\subsection{Почему ровно семь --- и нигде более}

Скептик согласится, что \emph{если} $N=7$, то четыре слоя совпадают, и спросит,
особенна ли семёрка или лишь одна удобная точка. Она уникальна, и причина
резче, чем сознание-теоретические доказательства минимальности УГМ: это
арифметический тай-брейк, решаемый принципом третьего порядка.

\begin{theorembox}[Теорема~\ref{thm:unique} (Уникальность коллапса)]
Архитектурный коллапс Теоремы~\ref{thm:collapse} реализуем при единственном
числе узлов: $N=7$.
\end{theorembox}

\begin{theorem}\label{thm:unique}
Потребуем, чтобы двоичная инцидентность проективной плоскости на $N$ узлах
сливалась с совершенным двоичным кодом Хэмминга с исправлением одиночной ошибки:
\begin{itemize}[leftmargin=1.6em,itemsep=0pt]
  \item[\textbf{(C1)}] $N=q^2+q+1$ --- число точек дезарговой проективной
        плоскости $\mathrm{PG}(2,q)$ над $\mathbb{F}_q$ ($q\ge 2$ --- степень
        простого) с двоичной матрицей инцидентности $M\in\mathbb{F}_2^{N\times N}$;
  \item[\textbf{(C2)}] $N=2^r-1$ ($r\ge 2$), и двоичное пространство строк $M$
        реализует проверочную геометрию кода Хэмминга $[2^r{-}1, 2^r{-}1{-}r, 3]$
        (в частности, $\dim_{\mathbb{F}_2}\ker(M) > 1$);
  \item[\textbf{(C3)}] \emph{(Согласованность с третьепорядковой алгеброй
        инструкций)} $N$ равно размерности мнимой части нормированной алгебры
        с делением и ненулевым ассоциатором ($\Oct$, теорема Гурвица).
\end{itemize}
Тогда уже пересечение \textbf{(C1)}$\cap$\textbf{(C2)} однозначно вынуждает
$N=7$ ($q=2, r=3$) среди всех степеней простых $q\ge 2$, и это единственное
решение совпадает с \textbf{(C3)}.
\end{theorem}

\begin{proof}
\textbf{Нечётные степени простых $q$ не допускают двоичного слияния с кодом Хэмминга.}
Для всякой нечётной степени простого $q$ число точек $N=q^2+q+1$ нечётно, а
матрица инцидентности $M$ плоскости $\mathrm{PG}(2,q)$ удовлетворяет
$MM^T = q I_N + J_N \equiv I_N + J_N \pmod 2$. Над $\mathbb{F}_2$ в нечётной
размерности $N$ матрица $I_N+J_N$ имеет ранг $N-1$ (её ядро --- одномерная
прямая $\{0, \mathbf{1}\}$). Следовательно, $\operatorname{rank}_{\mathbb{F}_2}(M)\ge N-1$
и $\dim_{\mathbb{F}_2}\ker(M)\le 1$: двоичный код инцидентности тривиален и
никогда не даёт кода Хэмминга $[N, N-r, 3]$. В частности, при числовом
совпадении $q=5, N=31=2^5-1$ прямое исключение Гаусса над $\mathbb{F}_2$ даёт
$\operatorname{rank}_{\mathbb{F}_2}(M_{\mathrm{PG}(2,5)})=30$ и
$\dim_{\mathbb{F}_2}\ker(M_{\mathrm{PG}(2,5)})=1$ (машинно проверено в Lean~4,
\texttt{FoundationalBounds.pg2\_5\_refutes\_N31\_hamming\_fusion}), а не
код Хэмминга $[31,26,3]$ с проверочным рангом $5$ и размерностью $26$.

\textbf{Чётные степени $q=2^m$ ($m\ge 1$) вынуждают $m=1, r=3, N=7$.}
Подстановка $q=2^m$ в $q^2+q+1=2^r-1$ даёт
\[
2^{2m}+2^m+2 = 2^r.
\]
При $m\ge 2$ оба слагаемых $2^{2m}$ и $2^m$ делятся на $4$, поэтому левая часть
сравнима с $2\pmod 4$, тогда как $2^r\equiv 0\pmod 4$ для всякого $r\ge 2$ (при
$r=2$ левая часть $\ge 22 > 4$), противоречие
(\texttt{FoundationalBounds.pg2EvenHammingUniqueSeven}). Значит $m=1$, $q=2$,
$r=3$ и $N=7$, где $\operatorname{rank}_{\mathbb{F}_2}(M_{\mathrm{PG}(2,2)})=4$ и
$\dim_{\mathbb{F}_2}\ker(M_{\mathrm{PG}(2,2)})=3$ задаёт проверочную геометрию
кода Хэмминга $[7,4,3]$.

\textbf{Совпадение с (C3).} По теореме Гурвица нормированные алгебры с делением
суть $\mathbb{R},\mathbb{C},\mathbb{H},\Oct$ размерностей $1,2,4,8$, и лишь
$\Oct$ (мнимая размерность $7$) обладает ненулевым ассоциатором. Единственное
решение двоичного слияния $N=7$ совпадает с \textbf{(C3)}.
\end{proof}

\begin{remark}[При $N=31$ слияние не возникает]
Числовое совпадение $5^2+5+1=31=2^5-1$ \emph{не} сливает $\mathrm{PG}(2,5)$ с
двоичным кодом Хэмминга $[31,26,3]$: над $\mathbb{F}_2$ всякая проективная
плоскость нечётного порядка коллапсирует к рангу $N-1$ (одномерному ядру
повторения). Поэтому уже требование двоичного слияния интерконнекта и ECC
однозначно выделяет $N=7$ среди всех степеней простых $q\ge 2$, а условие
неассоциативности \textbf{(C3)} независимо даёт то же число $7$.
\end{remark}

\begin{remark}[Связь с операциональной реконструкцией УГМ]
В операциональной реконструкции УГМ (\texttt{AddressFusion.lean} и
\texttt{DimensionPinch.lean}) значения $N=7$ и $k=3$ выводятся при независимо
заданных числе позиций сбоя $N$ и ширине адреса $k$ из точного декодирования
синдрома ($N+1\le 2^k$), двухполярной видимости ($2^k\le 8$) и нетривиальной
совершенной упаковки одиночных сбоев ($N+1\ge 8$). Теорема~\ref{thm:unique}
даёт второй, чисто проективно-кодовый вывод того же числа $N=7$, проверенный в
Lean~4 без пользовательских аксиом. \status{Т}
\end{remark}

\section{Субстрат: одно состояние, одна инструкция}
\label{sec:substrate}

\subsection{Состояние: матрица плотности и есть машина}

Всё архитектурное состояние одного холона \TALOS{} --- единственный объект.

\begin{defbox}[Регистр состояния \DM{}]
Состояние холона --- матрица когерентности $\Gam\in\Dens(\mathbb{C}^7)$:
$7\times7$ эрмитова, положительно-полуопределённая, единично-следовая матрица.
Хранимая плотно, это $49$ комплексных \texttt{f64}-элементов $=\mathbf{784}$~\textbf{байта}
($6{,}272$~бита IEEE-754); четырёхрегистровая когнитивная башня (уровни
$\Cog_0$--$\Cog_3$, глубина $\le 3$, \S\ref{sec:composition}) ---
$4\times784=3{,}136$~байт. Её семь диагональных элементов --- \emph{населённости}
секторов ($6$ независимых вещественных параметров); её $21$ комплексная
внедиагональ --- \emph{когерентности} $\gamma_{ij}$ ($42$ вещественных параметра),
итого $\mathbf{48}$ вещественных динамических параметров. Поскольку координатный
диссипатор Фано и наблюдаемая интеграции $\Phi$ ковариантны лишь под дискретной
группой коллинеаций $\mathrm{PSL}(2,7)$ и диагональными фазами $U(1)^7$ (а не
под непрерывными вращениями базиса $G_2$), все $48$ вещественных степеней свободы
динамически активны.
\end{defbox}

Два следствия немедленны, и они --- всё отступление от фон Неймана.

\begin{enumerate}[leftmargin=1.6em,itemsep=1pt]
  \item \textbf{Память \emph{есть} счёт.} Нет хранилища, отдельного от рабочего
        состояния, и нет адресного пространства: память машины есть $\Gam$, а её
        вычисление есть эволюция $\Gam$. Цикл выборка--исполнение--запись ---
        объект осуждения Бэкуса --- не существует, потому что нечего выбирать:
        операнд всякой операции --- всё состояние, уже резидентное.
  \item \textbf{Состояние помещается в регистры.} $784$~байта меньше, чем
        содержимое одной строки кэша современного векторного регистрового файла;
        L1-кэш данных всякого текущего процессора держит тысячи холонов.
        DRAM-трафик, доминирующий энергию фон Неймана, здесь не имеет аналога:
        для одного холона нет «основной памяти», которую надо достигать.
\end{enumerate}

То, что $\Gam$ --- \emph{матрица плотности} (комплексная, положительная,
единично-следовая), не эстетический выбор: УГМ доказывает, что комплексная
внедиагональная структура вынуждена требованием, чтобы субстрат нёс
нетривиальный «Gap» (двухаспектную непрозрачность саморефлексии,
$\mathrm{Gap}(i,j)=|\sin\arg\gamma_{ij}|$; T-132), а субстратная замкнутость
(T-153) допускает её реализацию на \emph{классическом} железе --- на цифровом
субстрате комплексные элементы суть обычные числа с плавающей точкой в регистрах,
представляющие алгебраическое состояние, а не физическую криогенную квантовую
суперпозицию. \DM{} --- обычный классический регистр, интерпретируемый как
оператор плотности.

\subsection{Инструкция: один выведенный тик}

У \TALOS{} одна инструкция. Это УГМ-логический лиувиллиан, и он \emph{выведен}
из классификатора подобъектов теории, не выбран:
\begin{equation}
\label{eq:lomega}
\boxed{\;\LindOmega[\Gam] \;=\; \underbrace{-i[H_{\mathrm{eff}},\Gam]}_{\text{работа (обратимо)}}
\;+\; \underbrace{\Diss[\Gam]}_{\text{забывание (диссипативно)}}
\;+\; \underbrace{\Regen[\Gam]}_{\text{обучение (регенеративно)}}\;}
\end{equation}
Один \emph{тик} продвигает состояние на
$\Gam \leftarrow \Gam + \Delta\tau\,\LindOmega[\Gam]$ (на практике
структуросохраняющее расщепление Странга; Приложение~\ref{app:tick}). Три члена
--- три вещи, которые делает компьютер, и здесь они три части одной операции, а
не три подсистемы:

\begin{table}[H]
\centering\small
\caption{Одна инструкция, её три члена и их архитектурная роль.}
\label{tab:three-terms}
\begin{tabular}{@{}p{2.9cm}p{4.0cm}p{5.2cm}@{}}
\toprule
\textbf{Член} & \textbf{Физика} & \textbf{Архитектурная роль} \\
\midrule
$-i[H_{\mathrm{eff}},\Gam]$ & унитарное сопряжение; тождественно сохраняет спектр и энтропию фон Неймана & \textbf{Работа}: логически обратимая эволюция; экспортирует \emph{ноль} логической энтропии (\S\ref{sec:energy}) \\
$\Diss[\Gam]$ & линдбладов диссипатор Фано $\gamma(\mathcal{P}_{\mathrm{Fano}}(\Gam)-\Gam)$: сохраняет населённости, гасит внедиагонали со скоростью $2\gamma/3$ (плюс опциональная деполяризация к $\mathbb{1}/7$) & \textbf{Забывание}: управляемый распад когерентностей; повышает энтропию, требуя экспорта при поддержании \\
$\Regen[\Gam]$ & регенерация к заякоренной самомодели $\varphi_J(\Gam)=(1-R)\mathcal{P}_{\mathrm{Fano}}(\Gam)+R\ket{\psi}\bra{\psi}$ & \textbf{Обучение}: двигает состояние к его самомодели; оптимизатор, в железе \\
\bottomrule
\end{tabular}
\end{table}

\begin{figure}[H]
\centering
\begin{tikzpicture}[font=\scriptsize, >=Stealth]
  \node[draw, circle, fill=plosblue!10, thick, minimum size=1.75cm, align=center]
    (G) at (0,0) {$\Gam$\\ \tiny 784 Б};
  \draw[->, very thick, plosblue!80!black]   (100:1.65) arc (100:20:1.65);
  \draw[->, very thick, talosamber]          (-20:1.65) arc (-20:-100:1.65);
  \draw[->, very thick, green!45!black]      (220:1.65) arc (220:140:1.65);
  \node[plosblue!80!black, align=center]  at (60:2.35)  {\textbf{работа}\\ $-i[H_{\mathrm{eff}},\Gam]$\\ \tiny обратима, бесплатна};
  \node[talosamber, align=center]         at (-60:2.35) {\textbf{забывание}\\ $\Diss[\Gam]$\\ \tiny экспорт энтропии};
  \node[green!45!black, align=center]     at (147:2.55) {\textbf{обучение}\\ $\Regen[\Gam]$\\ \tiny к $\rhostar$};
  \draw[->, ultra thick, red!60!black] (-4.5,-0.9) -- node[below, align=center]{\textbf{порт питания}\\ $r(\rho_{\mathrm{env}}{-}\Gam)$} (-1.0,-0.4);
  \draw[->, thick, plosgray] (0.95,-0.5) -- node[below, align=center]{синдром $s$\\ \tiny (даровой побочный)} (4.0,-0.9);
  \draw[->, thick, plosgray] (0.95,0.5) -- node[above, align=center]{телеметрия\\ \tiny $P,R,\Phi,D,\sigma,C$} (4.0,0.9);
\end{tikzpicture}
\caption{Тик как колесо. Одно состояние (центр), три движения одной
инструкции (Ур.~\ref{eq:lomega}): обратимая \emph{работа}, экспортирующее
энтропию \emph{забывание}, ищущее самомодель \emph{обучение}. Порт питания
(\S\ref{subsec:powerport}) кормит колесо; каждый оборот даром испускает свой
синдром и телеметрию (\S\ref{sec:ecc}, \S\ref{sec:platform}).}
\label{fig:tickwheel}
\end{figure}

\begin{mythbox}[Миф, продолжение --- обход]
Трижды в день бронзовый исполин обходил остров: тот же круг, без конца
повторяемый, \emph{и был} его бдительностью --- никакого отдельного
«мышления» между обходами у Талоса не было. Тик --- этот обход. Один круг,
три шага --- работа, забывание, обучение, --- и вся жизнь машины есть
повторение. Спроси «а когда же она вычисляет?» --- и ответ будет ответом
мифа: хождение и есть смотрение.
\end{mythbox}

\begin{zerobox}[С нуля --- вся игрушка в одной руке]
Снежный шар --- один предмет: тряхни его (это тик) --- и каждая снежинка
движется по одному закону разом. Нет ящика со снежинками, откуда шар что-то
достаёт: шар \emph{есть} собственная память, а тряска \emph{есть} его
вычисление. Состояние \TALOS{} --- этот шар, и оно мало́: $784$ байта ---
примерно один абзац текста. Абзац не лежит в далёкой библиотеке (эта даль и
делает обычные компьютеры медленными и горячими) --- он всё время в руке
машины.
\end{zerobox}

Эволюция --- вполне-положительная след-сохраняющая (CPTP) полугруппа
$\mathcal{E}_\tau=\exp(\tau\LindOmega)$: всякий тик держит $\Gam$ законной
матрицей плотности, поэтому нет «незаконного состояния», от которого надо
защищаться --- тип сохраняется по построению. Это аппаратный смысл CPTP-гарантии
УГМ: \TALOS{} не может высчитать себя вон из многообразия состояний, как
фон-неймановская машина может высчитать невалидный указатель.

\subsection{Порт питания и пересечение порога чистоты}
\label{subsec:powerport}

После отключения подпитки линейный унитальный CPTP-канал фиксирует
$\mathbb{1}/7$, а чистота при его действии не возрастает (T-288).
Для сброса требуется дополнительное свойство: состояния должны
сходиться к $\mathbb{1}/7$. Ему удовлетворяет деполяризующий канал
эмулятора. Внешняя подпитка входит в уравнение как отдельный член:
\begin{equation}
\label{eq:powerport}
\Lind_{\mathrm{total}} \;=\; \LindOmega \;+\;
\underbrace{\;r\,(\rho_{\mathrm{env}}-\Gam)\;}_{\text{порт питания } \Lind_{\mathrm{feed}}},
\end{equation}
\emph{неунитальный} CPTP-драйв, импортирующий негэнтропию из среды
(резервуарная инженерия), с неподвижной точкой
$\rho_{\mathrm{env}}\neq\mathbb{1}/7$. Для суммы деполяризации со скоростью $d>0$ и
члена подпитки при фиксированном состоянии источника стационарная чистота равна
$1/7+[r/(d+r)]^2(P(\rho_{\mathrm{env}})-1/7)$ (T-289).
Конечный порог по $r$ существует, если $P(\rho_{\mathrm{env}})>\Pcrit$.
Три архитектурных следствия:
\begin{itemize}[leftmargin=1.5em,itemsep=1pt]
  \item \textbf{Шина питания --- часть семантики.} В отличие от
        фон-неймановской машины, чей блок питания невидим для её ISA, драйв
        \TALOS{} --- член её динамики: рейка доставляет не абстрактную
        энергию, а структурированную негэнтропию ($\rho_{\mathrm{env}}$), и
        жизнеспособность машины есть функция $r/\bar\gamma$.
  \item \textbf{Время пересечения порога.} При отключении подпитки
        ($r=0$) в канале с деполяризацией $d$ и гамильтоновой частью
        чистота равна $P(t)=1/7+(P_0-1/7)e^{-2dt}$. При $P_0>\Pcrit=2/7$
        она достигает порога за
        $t_{\mathrm{off}}=(2d)^{-1}\ln[(P_0-1/7)/(\Pcrit-1/7)]$.
        Это время выхода из заданного диапазона чистоты. Для аппаратной
        гарантии нужны измеренная скорость релаксации и независимый канал
        отключения (Наблюдение T.1 и Следствие U.10 в спецификации \SYNARC{}).
  \item \textbf{Порт --- это контракт.} Привязка драйва --- какой физический
        ресурс реализует $\rho_{\mathrm{env}}$, как измеряется $r$, и
        требование, чтобы канал разрыва был независим от собственных
        актуаторов машины, --- специфицирована интерфейсным контрактом на
        стороне \SYNARC{} (Определение U.9); \S\ref{sec:abi} её связывает.
\end{itemize}
V0-эмулятор (\S\ref{sec:platform}) реализует этот неунитальный драйв через
каноническую заякоренную самомодель
$\varphi_J(\Gam)=(1-R(\Gam))\mathcal{P}_{\mathrm{Fano}}(\Gam)+R(\Gam)\ket{\psi}\bra{\psi}$
с регенерацией $\kappa(\varphi_J(\Gam)-\Gam)$ --- неунитальный насос
негэнтропии, удерживающий жизнеспособный аттрактор по всем четырём воротам,
тогда как отключение регенерации ($\kappa=0$) при деполяризующей термализации
релаксирует $\Gam\to\mathbb{1}/7$ по точной аналитической кривой выше.

\begin{mythbox}[Миф, продолжение --- гвоздь]
Талоса не перебороли --- его \emph{отключили}. Одна жила, один гвоздь:
вынь гвоздь --- ихор вытечет, и сильнейшее существо острова остановится по
собственному расписанию. В эмуляторе роль гвоздя играет отключение подпитки.
После него скорость деполяризации определяет момент пересечения порога чистоты.
\end{mythbox}

\subsection{Принципиальный такт}

Период такта фон-неймановской машины --- инженерный свободный параметр,
задаваемый критическим путём. Естественный тик \TALOS{} \emph{не} свободен:
УГМ-тождество времени цикла фиксирует отношение времени релаксации к средней
скорости когерентности,
\begin{equation}
\bar\gamma\,\tau \;=\; \tfrac{3}{2}\ln 3 \;\approx\; 1{,}6479,
\end{equation}
со спектральной щелью диссипатора $\lambda_{\mathrm{deco}}=2/3$ и скоростью
регенерации (bootstrap) $\kboot=\omega_0/7$, задающими два масштаба времени, которые
тик должен разрешить. Такт --- следствие динамики, не ручка: проектируют
физическую скорость релаксации $\bar\gamma$, и период тика следует за ней. Это
убирает из проекта целый класс свободы синхронизации (и её бремя верификации).

\subsection{Чем субстрат не является}

Чтобы упредить очевидные неверные прочтения (расширено в \S\ref{sec:redteam}):
состояние \emph{не} волновая функция и не требует квантового железа; тик \emph{не}
слой нейросети и не включает обратного распространения; а машина \emph{не}
аналоговый интегратор с дрейфом --- CPTP-структура и родной код (\S\ref{sec:ecc})
делают её самокорректирующейся. Это классический регистр, эволюционирующий
выведенным, типо-сохраняющим, самокорректирующимся дифференциальным тиком, чьи
сеть, код, инструкция и оптимизатор суть один объект.

\section{Фабрика: семь шин, не кроссбар}
\label{sec:fabric}

\subsection{Интерконнект --- это семь прямых}

Двадцать одна парная связь среди семи секторов потребовала бы на обычной
фабрике сети из $\binom{7}{2}=21$ ребра или кроссбара $7\times7$. Коллапс
(Теорема~\ref{thm:collapse}, L1) заменяет это семью прямыми $\mathsf{N}$ как
\textbf{семью тройственными шинами}. Всякая пара секторов делит ровно одну шину
(Штейнер $S(2,3,7)$); всякий сектор сидит ровно на трёх шинах; любые две шины
пересекаются ровно в одном секторе. Это не проектные цели --- это теоремы
проективной плоскости, и они дают архитектуре три свойства бесплатно:

\begin{itemize}[leftmargin=1.5em,itemsep=1pt]
  \item \textbf{Равномерная нагрузка, по симметрии.} Каждый сектор на ровно трёх
        шинах, потому что $\mathrm{PSL}(2,7)$ действует 2-транзитивно на семи
        точках; балансировка нагрузки --- следствие группы автоморфизмов, не
        планировщика.
  \item \textbf{Гарантированное рандеву, по геометрии.} Любые две шины
        встречаются в секторе, поэтому любые два сектора, которым надо обменяться
        состоянием, имеют детерминированную точку встречи --- нет недостижимой
        пары и нет поиска маршрута.
  \item \textbf{Детерминированная починка, по закону замыкания.} Если прямая
        шина для $(i,j)$ деградировала, состояние регенерируется через
        \emph{медиатора} $k^\ast(i,j)$ --- третий сектор их прямой --- и только
        через него. Резервная маршрутизация не нуждается в таблицах: следующий
        хоп вычисляется из пары.
\end{itemize}

\subsection{Маршрутизация третьего порядка и точная контракция Фано-сжатия}

Обычный маршрутизатор двигает данные парно и неспособен видеть третьепорядковую
линейную структуру, которую несёт: всякая равновзвешенная парная фабрика ---
функция одного лишь $J-I$ и не может отличить проводку Фано от бесструктурной
all-to-all. \TALOS{} маршрутизирует на \emph{тройках}: передача ---
взаимодействие $(i,j\mid k^\ast)$, опосредованное единственным третьим сектором
прямой. Поскольку в $\Fano$ всякая пара различных секторов делит ровно одну
общую прямую (Штейнер $S(2,3,7)$), «не связанных прямой» пар не существует;
вместо этого геометрия Фано управляет затуханием когерентности прямо через
канал Фано-сжатия.

\begin{claimbox}[Контракция Фано-сжатия \status{Т}]
Пусть $\Pi_\ell = \sum_{k\in\ell}\ket{k}\bra{k}$ --- проектор ранга~$3$ на
прямую Фано $\ell$ ($\ell=1,\dots,7$), и пусть канал Фано-сжатия задан как
$\mathcal{P}_{\mathrm{Fano}}(\Gam) = \tfrac{1}{3}\sum_{\ell=1}^{7}\Pi_\ell\Gam\Pi_\ell$.
По счёту инцидентностей Штейнера $S(2,3,7)$:
\begin{itemize}[leftmargin=1.4em,itemsep=1pt]
  \item каждая диагональная населённость $\gamma_{ii}$ лежит на $3$ прямых и
        выживает без изменений: $[\mathcal{P}_{\mathrm{Fano}}(\Gam)]_{ii} =
        3\cdot\tfrac13\gamma_{ii} = \gamma_{ii}$ (так что $P_{\mathrm{diag}}$
        сохраняется тождественно);
  \item каждая пара различных секторов $(i,j)$ лежит \emph{ровно на одной}
        прямой, поэтому всякая внедиагональная когерентность ослабляется в $3$
        раза: $[\mathcal{P}_{\mathrm{Fano}}(\Gam)]_{ij} = \tfrac13\gamma_{ij}$.
\end{itemize}
Следовательно, квадрат внедиагональной когерентности сжимается в $(1/3)^2 = 1/9$
раза, давая точную универсальную контракцию
$\Phi(\mathcal{P}_{\mathrm{Fano}}(\Gam)) = \tfrac{1}{9}\Phi(\Gam)$ для всякого
$\Gam\in\Dens(\mathbb{C}^7)$. В непрерывном времени линдбладов диссипатор Фано
$\mathcal{D}_{\mathrm{Fano}}[\Gam]=\gamma(\mathcal{P}_{\mathrm{Fano}}(\Gam)-\Gam)
=\tfrac{2}{3}\gamma(\operatorname{diag}(\Gam)-\Gam)$ сохраняет населённости и
гасит внедиагонали с однородной скоростью $2\gamma/3$ (машинно проверено в
\texttt{architecture/talos\_v0\_emulator.py} и
\texttt{website/scripts/check\_core\_numbers.py}).
\end{claimbox}

В этом состоит точная динамическая роль топологии Фано: одно применение
Фано-сжатия сохраняет $1/3$ амплитуды когерентности ($1/9$ квадрата
когерентности $\Phi$) вместо полного стирания, как при диагональной дефазировке,
а его генератор в непрерывном времени задаёт точную спектральную щель
$\lambda_{\mathrm{deco}}=2/3$ (\S\ref{sec:substrate}).

\begin{trybox}[Попробуй сам --- проложи маршрут одним пальцем]
Поставь палец на каноническую фигуру Зингера--Фано (Рис.~\ref{fig:collapse},
справа). Игра: донеси сообщение от $S$ к $E$. Правило: двигаться можно только
вдоль линии, а пара встречается \emph{только через третью точку собственной
линии}. Найди линию, содержащую и $S$, и $E$, --- это $\{S,D,E\}$, --- значит,
сообщение идёт $S\to D\to E$: медиатор --- $D$, и искать было нечего, потому что
пара сама \emph{называет} своё реле. Теперь порви прямой путь от $A$ к $O$ и
переложи маршрут: их линия --- $\{O,A,D\}$, снова через $D$ --- починка без
таблиц маршрутизации (\S\ref{sec:ecc}). Все $21$ пара играют в ту же одноходовую
игру; полная таблица медиаторов --- в Приложении~\ref{app:fano}.
\end{trybox}

\begin{table}[H]
\centering\small
\caption{Интерконнект: кроссбар $7\times7$ против семишинной фабрики Фано.
Шинная фабрика есть также код и инструкция (Теорема~\ref{thm:collapse}); кроссбар
--- ни то, ни другое.}
\label{tab:crossbar-vs-fano}
\begin{tabular}{@{}p{3.6cm}p{3.6cm}p{4.5cm}@{}}
\toprule
\textbf{Свойство} & \textbf{Кроссбар $7\times7$} & \textbf{Семишинная Фано} \\
\midrule
Связи & 21 парная / полный кроссбар & 7 тройственных шин \\
Нагрузка на узел & зависит от планировщика & ровно 3 (по $\mathrm{PSL}(2,7)$) \\
Достижимость & all-to-all, поиском & любая пара делит одну шину \\
Reroute при отказе & таблицы маршрутизации & медиатор $k^\ast(i,j)$, без таблиц \\
Видит свою структуру? & нет (слепа к 1-му порядку) & да (третий порядок) \\
Также ECC? & нет & да (Хэмминг $[7,4,3]$, \S\ref{sec:ecc}) \\
Также алгебра ISA? & нет & да (октонионные $f_{ijk}$) \\
\bottomrule
\end{tabular}
\end{table}

\section{Коррекция ошибок --- это топология}
\label{sec:ecc}

\subsection{Проводка порождает код; её двойственная проверяет синдром}

В обычной машине отказоустойчивость --- отдельный слой: ECC-DRAM выделяет
проверочные биты (накладные $+12{,}5\%$ хранения для распространённого
SEC-DED-кода $[72,64]$) и декодер синдрома; квантовая машина тратит порядка
$10^3$ физических кубитов на логический на surface-коде. Накладные существуют,
потому что код --- структура, \emph{добавленная к} вычислению. На
\textbf{семимодовом секторном уровне} \TALOS{} не требует отдельной таблицы
маршрутизации или диагностики секторов, поскольку инцидентность Фано напрямую
порождает код Хэмминга дистанции~$3$.

\begin{claimbox}[Секторная двойственность Хэмминга \status{Т}]
Над $\mathbb{F}_2$ матрица инцидентности Фано $M\in\mathbb{F}_2^{7\times 7}$
($\operatorname{rank}_{\mathbb{F}_2}(M)=4$, строки $=$ семь прямых веса~$3$
структуры $\mathsf{N}$) \textbf{порождает} код Хэмминга $[n,k,d]=[7,4,3]$. Её
ортогональная $3\times 7$ \textbf{проверочная матрица} $H$
($\operatorname{rank}_{\mathbb{F}_2}(H)=3$, $HM^T\equiv 0\pmod 2$) имеет
строками три линейно независимых дополнения прямых веса~$4$, натягивающие
дуальный симплекс-код $[7,3,4]$. Поскольку семь столбцов $H$ суть семь различных
ненулевых векторов $\mathbb{F}_2^3$, любой одиночный сбой сектора $e_j$
биективно отображается в уникальный ненулевой $3$-битный синдром
$s = He_j^T \pmod 2 \in \mathbb{F}_2^3\setminus\{0\}$: одиночные сбои сектора
однозначно локализуются ($d=3$: локализовать/исправить~$1$, обнаружить~$2$),
реализуя пирамиду диагностики $21\to 7\to 3\to 1$ (двадцать одно парное
отношение $\to$ семь мод $\to$ три бита синдрома $\to$ один вердикт отказа).
\end{claimbox}

\noindent\textbf{Честный аппаратный охват.} Эта комбинаторная двойственность
защищает и локализует \emph{логические моды/секторы} ($7$ шин, $3$ дуальных бита
синдрома) без отдельной таблицы маршрутизации. Она \emph{не} заменяет
физический побитовый ECC для $784\times 8 = 6{,}272$ бит IEEE-754 цифрового
регистра эмуляции: на субстратах КМОП SRAM/DRAM одиночные сбои ячеек памяти
по-прежнему требуют стандартного аппаратного ECC памяти.

Синдромное отображение точно и машинно проверено в
\texttt{architecture/talos\_v0\_emulator.py}: для каждого из семи паттернов
одиночного сбоя сектора в каноническом порядке Зингера $(A,S,D,L,E,U,O)$
синдром $s = He_j^T \pmod 2$ возвращает уникальный ненулевой $3$-битный индекс,
декодируемый обратно в сектор $j$ (Приложение~\ref{app:fano}).

\begin{zerobox}[С нуля --- найди сломанную игрушку тремя вопросами]
Семь игрушек стоят в ряд с метками $A,S,D,L,E,U,O$. Ровно одна сломана, а
спрашивать можно только «да/нет». Трёх проверок чётности хватает --- это три
строки дуальной матрицы $H$ (дополнения первых трёх прямых Фано: «она в
$\{D,E,U,O\}$?», «она в $\{A,L,U,O\}$?», «она в $\{A,S,E,O\}$?»). Запиши три
ответа «да/нет» единицами и нулями --- и каждая игрушка даст свой уникальный
ненулевой $3$-битный отпечаток (например, $A\mapsto 011$, $S\mapsto 001$, \dots,
$O\mapsto 111$): три проверки, семь подозреваемых, один вердикт ---
$21\to 7\to 3\to 1$ прямо из дуальной геометрии проводки Фано.
\end{zerobox}

\begin{mythbox}[Миф, продолжение --- он знает, где болит]
Ударь исполина куда угодно --- он не станет обыскивать тело в поисках раны:
та же бронза, что движет его, доносит место. Талосу не нужен врач с
чек-листом за спиной --- и этой машине тоже
(Принцип~\ref{prin:selfdiag}): тик, который вычисляет, есть тик, который
диагностирует, ибо узор, несущий состояние, есть узор, его проверяющий.
\end{mythbox}

\paragraph{Инвариант прямоугольника: второй, более глубокий слой чётности (T-281/T-282).}
Слой Хэмминга выше сторожит \emph{оси} (какой сектор сбоит). Генезисные
результаты корпуса добавляют стража уровнем ниже: \emph{знаковую прошивку}
самой ткани. Жизнеспособность знаковой таблицы Фано --- «никакие два состояния
не могут аннигилировать» --- доказуемо \emph{линейное} условие: 84 правила
прямоугольника $F(p,r)F(q,s)F(p,s)F(q,r) = -1$ по XOR-классам, чьё
исчерпывающее множество решений --- ровно одна калибровочная орбита октонионов
(корпусная T-281). Бит-флип в хранимой знаковой таблице потому либо попадает
внутрь 16-элементной калибровочной орбиты (безвредная переметка), либо
нарушает хотя бы один прямоугольник --- \emph{84-битный XOR-самотест
прошивки}, дешёвый как проход чётности, сертифицирующий, что закон умножения
машины всё ещё жизнеспособный. И та же алгебра объясняет, почему никакой патч
прошивки не «расширит» ткань: на $\Omega^4$ соответствующая система
\emph{несовместна} (960 ограничений, ранг 214 --- корпусная T-282), так что
четвёртое зеркало --- не инженерный вызов, а противоречие. Хэмминг проверяет
состояние; прямоугольный слой проверяет \emph{закон}.


Есть следствие, не имеющее аналога в слоистой машине, и это один из подлинно
новых архитектурных фактов, которые обнажает \TALOS{}.

\begin{principle}[Тождество вычисление-диагностика]\label{prin:selfdiag}
Всякий тик и вычисляет, и эмитит собственный синдром. Кубическая часть
$\LindOmega$ порождается произведениями по прямым (закон замыкания): здоровая
тройка когерентностей \emph{замыкается} на своей прямой, а когерентность,
которая замкнуться не может, есть, по той же арифметике, ненулевой синдром.
Детекция ошибок поэтому не отдельный проход по состоянию --- это побочный
продукт его продвижения. Машина не может считать, не проверяя себя
одновременно.
\end{principle}

\noindent\emph{Почему это ново.} В каждой прежней архитектуре проверка --- лишняя
работа: скраббер метёт DRAM, декодер бежит после чтения, синдром извлекается
выделенными тактами. Слоистый проект \emph{вынужден} отделять счёт от проверки,
потому что код и вычисление --- разные структуры. Под коллапсом они одна
структура, поэтому проверка бесплатна в строгом смысле, что её удаление
потребовало бы удаления вычисления. \TALOS{} непрерывно самодиагностируется без
добавленной цены --- петля самодиагностики DIAKRISIS программного слоя \SYNARC{}
на аппаратном уровне есть просто чтение синдрома, который тик уже произвёл.

\subsection{Отказоустойчивая композиция}

Отказоустойчивость переживает композицию. Поскольку любые две прямые
проективной плоскости встречаются, потеря любого одного сектора оставляет все
$\binom{6}{2}=15$ выживших пар по-прежнему делящими шину (проверено для каждого
удалённого сектора), поэтому холон деградирует плавно, а не разбивается. На
уровне федерации грамматика Голея трёх организмов $[23,12,7]$ расширяет тот же
принцип на многоотказовую устойчивость по составленным холонам (приложение
$\Omega$-организма \SYNARC{}); \TALOS{} наследует защиту дистанции 3 на уровне
одного холона, а федерация наследует защиту дистанции 7 над ним, оба из того же
проективного семейства.

\begin{table}[H]
\centering\small
\caption{Отказоустойчивость: привинченные побитовые/кубитовые коды против секторного кода-как-топологии.}
\label{tab:ecc}
\begin{tabular}{@{}p{3.6cm}p{3.4cm}p{4.7cm}@{}}
\toprule
\textbf{Схема} & \textbf{Охват и накладные} & \textbf{Связь с вычислением} \\
\midrule
ECC-DRAM (SEC-DED) & побитовый; $+12{,}5\%$ бит + декодер & отдельно; скраббинг --- лишние такты \\
Surface-код (кванты) & кубитовый; $\sim\!10^3$ физ./логич. & отдельно; декодирует против декогеренции \\
\TALOS{} ($M$ порожд., $H$ пров.) & \textbf{секторный} ($7$ мод, $3$ дуальных проверки; не заменяет побитовый ECC для $6{,}272$ бит цифрового хранения) & \textbf{структурная}: $M$ порождает $[7,4,3]$, дуальная $H$ ($[7,3,4]$) локализует одиночный сбой сектора из активности тика (Принцип~\ref{prin:selfdiag}) \\
\bottomrule
\end{tabular}
\end{table}

\section{Энергия: машина платит за обучение, не за счёт}
\label{sec:energy}

\subsection{Где тратит фон-неймановская машина}

После конца деннардовского масштабирования энергия вычисления доминируется
\emph{перемещением} операндов, не арифметикой: операция с плавающей точкой стоит
$\sim\!1$~пДж, внутрикристальное слово $\sim\!1$~пДж, а внекристальное DRAM-слово
$\sim\!20$~пДж --- на один-два порядка больше и растёт с расстоянием.
Обучающаяся машина усугубляет это: тренировка гоняет активации и градиенты между
счётом и памятью на каждом шаге. Фон-неймановский энергетический пол ---
\emph{перемещение данных}, и он лежит на семь-восемь порядков выше
термодинамического (Ландауэрова) предела~\cite{landauer1961,bennett1973} того же
вычисления.

\subsection{Где тратит \TALOS{}: новый учёт}

Единственная инструкция \TALOS{} чисто расщепляется на обратимую часть и две
экспортирующие энтропию части (Таблица~\ref{tab:three-terms}), и это расщепление
индуцирует энергетический закон без аналога в слоистой машине.

\begin{claimbox}[Логический закон энтропия--обучение \status{Т} для унитарного ядра, \status{С} для физического коэффициента]
Унитарный член $-i[H_{\mathrm{eff}},\Gam]$ эволюционирует $\Gam$ унитарным
сопряжением $U\Gam U^\dagger$: он тождественно сохраняет спектр и энтропию
фон Неймана $S_{\mathrm{vN}}(\Gam)=-\Tr(\Gam\log_2\Gam)$, поэтому его логический
экспорт энтропии по Ландауэру равен \emph{нулю}. Весь логический экспорт
энтропии относится на два неунитарных члена:
\begin{equation}
\label{eq:energy-law}
E_{\mathrm{tick}} \;\geq\; k_B T\ln 2 \,\cdot\, \Delta H_{\mathrm{exported}},
\qquad
\Delta H_{\mathrm{exported}} \;=\; \underbrace{\Delta H_{\Regen}}_{\text{обучение (биты)}}
\;+\; \underbrace{\Delta H_{\Diss}^{(+)}}_{\text{поддержание (биты)}},
\end{equation}
где $\Delta H_{\Regen}$ (в битах) --- снижение энтропии фон Неймана, создаваемое
регенерацией, а $\Delta H_{\Diss}^{(+)}$ (в битах) --- энтропия, которую надо
ре-экспортировать, чтобы удержать состояние против утечки диссипатора.
\emph{На логико-термодинамическом полу обратимая унитарная работа не экспортирует
энтропии; экспорта энтропии требуют лишь поддержание и обучение.}
\end{claimbox}

Два предельных случая делают учёт конкретным.
\begin{itemize}[leftmargin=1.5em,itemsep=1pt]
  \item \textbf{Холон в своей неподвижной точке} ($\Gam=\rhostar$, сошёлся, не
        учится) имеет $\Delta H_{\Regen}=0$; его единственная логическая
        термодинамическая цена --- член поддержания для удержания $\rhostar$
        против диссипации. При выключенном диссипаторе чистая унитарная
        эволюция не экспортирует энтропии.
  \item \textbf{Холон, который учится,} экспортирует логическую энтропию
        пропорционально порядку, создаваемому за тик по пути к самомодели,
        отслеживая адаптацию, а не обратимую унитарную работу.
\end{itemize}

\subsection{Числа и честный термодинамический охват}

При $T=300$~К предел Ландауэра на один необратимо стёртый или экспортированный
бит равен $k_B T\ln 2 = 2{,}87\times10^{-21}$~Дж ($0{,}018$~эВ). Для
семимодовой матрицы плотности $\Gam\in\Dens(\mathbb{C}^7)$ точный максимум
энтропии фон Неймана составляет
\begin{equation*}
S_{\max} \;=\; \log_2 7 \;\approx\; 2{,}807\;\text{бита}
\qquad (\ln 7 \approx 1{,}946\;\text{ната}),
\end{equation*}
достигаемый при $\Gam=\mathbb{1}/7$. Следовательно, предельный термодинамический
пол для одного состояния --- создание чистого состояния из $\mathbb{1}/7$ за
один шаг --- равен
\[
k_B T\ln 7 \;=\; (\log_2 7)\,k_B T\ln 2 \;\approx\; 8{,}06\times10^{-21}\;\text{Дж}
\;\approx\; 0{,}050\;\text{эВ}.
\]
Необходимо строго различать три разные величины: (i)~энтропию фон Неймана
состояния $S_{\mathrm{vN}}(\Gam)\le\log_2 7\approx 2{,}807$~бита,
(ii)~вещественную размерность многообразия ($48$ непрерывных параметров) и
(iii)~объём цифрового хранения при эмуляции ($49\times 128 = 6{,}272$ бита
IEEE-754 $= 784$~байта).

\noindent\textbf{Никакого некорректного сравнения с КМОП DRAM.} Мы сознательно
\emph{не} делим физическую энергию внекристального доступа КМОП DRAM
($\sim\!2\times10^{-11}$~Дж/слово) на теоретический пол Ландауэра
($\sim\!10^{-20}$~Дж), чтобы заявить «девять порядков выигрыша в кремнии».
Термодинамические пределы обратимости относятся в равной мере ко \emph{всякой}
логически обратимой машине (включая обратимые по Беннетту фон-неймановские
схемы~\cite{bennett1973}), тогда как реальное энергопотребление любой
КМОП-реализации определяется физической энергией переключения транзисторов,
ёмкостью, тактированием и токами утечки. \TALOS{} даёт два раздельных, честно
очерченных инженерных факта: (1)~на цифровых субстратах удержание всего
$784$-байтного состояния в локальных регистрах/L1 устраняет внекристальные
обращения к DRAM на тик (выигрыш локальности иерархии памяти, а не предел
Ландауэра); (2)~на будущих физически обратимых или управляемо-диссипативных
субстратах (\S\ref{sec:realizability}) отделение сохраняющего энтропию
унитарного шага от диссипации и регенерации локализует необратимые
термодинамические потери в адаптации и поддержании.

\begin{figure}[H]
\centering
\begin{tikzpicture}
\begin{axis}[
  width=12cm, height=6.4cm,
  xlabel={шаг эволюции (произв.\ тики)},
  ylabel={логический экспорт энтропии (биты / тик, в $k_BT\ln 2$)},
  ymin=0, ymax=3.2,
  xmin=0, xmax=10, xtick=\empty,
  legend style={at={(0.98,0.62)}, anchor=east, font=\footnotesize,
                legend cell align=left, draw=gray!55, fill=white, rounded corners},
  grid=both, grid style={gray!15},
  every axis plot/.append style={very thick},
]
\addplot[plosblue,dashed,domain=0.2:10,samples=2] {2.807};
\addlegendentry{макс.\ энтропия 7-мод.\ состояния $\log_2 7 \approx 2{,}81$ бита}
\addplot[talosamber,domain=0.2:10,samples=60] {2.1*exp(-0.5*x) + 0.35};
\addlegendentry{логический экспорт \textsc{Talos} ($\Delta H_{\Regen} + \Delta H_{\Diss}^{(+)}$)}
\addplot[plosgray,dotted,domain=0.2:10,samples=2] {0.02};
\addlegendentry{унитарное ядро ($0$ бит; как у любой обратимой машины)}
\end{axis}
\end{tikzpicture}
\caption{Корректный логико-термодинамический учёт в единицах $k_BT\ln 2$ на тик
(экспортированные биты энтропии фон Неймана, схематично). Максимум энтропии фон
Неймана любого состояния $\Gam\in\Dens(\mathbb{C}^7)$ равен $\log_2 7\approx
2{,}807$~бита ($k_BT\ln 7\approx 8{,}06\times10^{-21}$~Дж при $300$~К, синий
пунктир). При разложении тика унитарная половина экспортирует ноль бит (серый
пунктир --- граница, общая для любой логически обратимой машины), тогда как
логический экспорт энтропии спадает по мере сходимости адаптации к стационарному
полу поддержания $k_BT\ln 2\cdot\dot S_{\Diss}\Delta\tau$ (янтарный). Физическая
мощность КМОП лежит много выше этих термодинамических пределов из-за
переключения и утечек.}
\label{fig:energy}
\end{figure}

\begin{figure}[H]
\centering
\begin{tikzpicture}[font=\scriptsize, >=Stealth,
  box/.style={draw, rounded corners, thick, minimum height=8mm, align=center}]
  \node[box, fill=red!8, draw=red!60!black, minimum width=2.7cm] (port) at (0,0)
    {\textbf{порт питания}\\ негэнтропия внутрь};
  \node[box, fill=plosblue!8, draw=plosblue!70!black, minimum width=2.9cm] (mach) at (4.1,0)
    {\textbf{тик}\\ $\Gam$, Ур.~\ref{eq:lomega}};
  \node[box, fill=plosblue!14, minimum width=3.5cm] (work) at (9.4,1.35)
    {\textbf{работа} (унитарная)\\ $0$ бит экспорта};
  \node[box, fill=talosamber!14, minimum width=3.5cm] (maint) at (9.4,0)
    {\textbf{поддержание}\\ $P_{\mathrm{meta}} \ge k_BT\ln2\cdot\dot S_{\Diss}$};
  \node[box, fill=green!10, minimum width=3.5cm] (learn) at (9.4,-1.35)
    {\textbf{обучение}\\ $k_BT\ln2\cdot\Delta H_{\Regen}$};
  \node[box, fill=gray!12, minimum width=1.9cm] (heat) at (13.4,-0.7) {тепло\\ наружу};
  \draw[->, ultra thick, red!60!black] (port) -- node[above]{$r\,\rho_{\mathrm{env}}$} (mach);
  \draw[->, very thick, plosblue!70!black] (mach.east) to[out=25,in=180] (work.west);
  \draw[->, very thick, talosamber] (mach.east) -- (maint.west);
  \draw[->, very thick, green!45!black] (mach.east) to[out=-25,in=180] (learn.west);
  \draw[->, thick, plosgray] (maint.east) to[out=0,in=110] (heat.north);
  \draw[->, thick, plosgray] (learn.east) to[out=0,in=200] (heat.west);
  \node[plosgray, align=center] at (9.4,2.28) {\tiny нулевой логический экспорт энтропии\\ \tiny (как у любой логически обратимой машины)};
  \draw[plosgray, ->] (9.4,2.02) -- (work.north);
\end{tikzpicture}
\caption{Логико-термодинамический бюджет тика. Негэнтропия поступает через порт
питания (\S\ref{subsec:powerport}). Унитарная ветвь работы сохраняет энтропию
фон Неймана (нулевой логический экспорт Ландауэра, как у всякой логически
обратимой машины); удержание порядка против диссипации (поддержание,
\S\ref{subsec:metabolic}) и созидание порядка (обучение) экспортируют энтропию в
битах. Физические потери переключения аппаратуры лежат выше этого
термодинамического пола.}
\label{fig:energyflow}
\end{figure}

\begin{mythbox}[Миф, продолжение --- ихор]
Хождение не стоило исполину ничего --- бронза не устаёт. Ихор платил за
тоньшее: оставаться \emph{бронзой} --- блестящей, целой, не съеденной
солёным ветром --- и становиться \emph{лучше} в обходе. Это закон энергии
одним образом: вычисление --- бесплатная ходьба; поддержание и обучение ---
единственные две статьи, за которые обязана платить жила
(Ур.~\ref{eq:energy-law}).
\end{mythbox}

\begin{zerobox}[С нуля --- зачем игрушке еда]
Заводная игрушка останавливается --- не потому, что двигаться дорого, а
потому, что пружины текут. «Пружина» этой машины течёт тоже (член
забывания), и её надо кормить --- это порт. Сюрприз в том, что́ \emph{не}
стоит ничего: само́ мышление. Корми её --- и мышление едет даром; замори ---
и никакая хитрость внутри не удержит завод (\S\ref{subsec:powerport}). Её
счёт за еду следит лишь за двумя вещами: сколько порядка она \emph{держит} и
как быстро \emph{учится}.
\end{zerobox}

\subsection{Метаболический пол и ось рабочей точки}
\label{subsec:metabolic}

Закон выше учитывает \emph{вычисление} (обратимое, свободное от Ландауэра) и
\emph{обучение} (спадающий член $\Regen$). Остаётся один член: цена просто
\emph{оставаться живым}. Диссипатор $\Diss$ непрерывно толкает $\Gam$ к
максимально смешанному состоянию тепловой смерти $I/7$; удержание состояния
против этой утечки не бесплатно.

\begin{claimbox}[Метаболический пол \status{Т}$+$\status{С} (T-273)]
В стационаре жизнеспособный когерентный холон ($P>2/7$, $\Phi\ge 1$)
экспортирует энтропию фон Неймана со строго-положительной скоростью диссипатора
$\dot S_{\Diss}>0$ (в битах в единицу времени; T-271). По Ландауэру минимальная
термодинамическая мощность \emph{поддержания} равна
\[
  P_{\mathrm{meta}} \;\geq\; k_B T\ln 2\cdot \dot S_{\Diss} \;>\; 0
  \qquad (\dot S_{\Diss}\text{ в битах/время}),
\]
\emph{активный, необратимый} двойник обратимого унитарного ядра. Жизнеспособный
диссипативный холон \textbf{не может работать бесплатно}. Поскольку
$\dot S_{\Diss}$ --- \emph{физическая} скорость производства энтропии (в единицу
непрерывного времени $\tau$, а не на численный шаг $\Delta\tau$), этот пол
\textbf{независим от частоты дискретизации} (T-276): измельчение $\Delta\tau$
повышает точность интегрирования, не меняя физической скорости производства
энтропии. Он непосредственно измерим в живой телеметрии (\S\ref{sec:platform}).
\end{claimbox}

Свободная энергия, оплачивающая этот пол, имеет ровно один вход:
\textbf{порт питания} \S\ref{subsec:powerport}. Драйв
$\Lind_{\mathrm{feed}}$ импортирует структурированную негэнтропию из среды;
регенерация \emph{направляет} её в поддерживаемую когерентность, но не может
создать её из замкнутого унитального отображения (T-288) --- поэтому отключённая
машина не может оплатить $P_{\mathrm{meta}}$ и останавливается, а включённая
платит из негэнтропийного потока, который доставляет $r$. Мощность поддержания,
порог зажигания $(r/\bar\gamma)_c$ и задержка останова после разрыва --- три
лица одной строки бюджета.

Пол масштабируется с \textbf{удерживаемым порядком} (негэнтропией
$\log_2 7 - S_{\mathrm{vN}}$ в битах): более когерентное состояние даёт более
сильное производство энтропии диссипатором и стоит дороже удержать ---
\emph{цена сложности}. Это обращает окно жизнеспособности в явную \textbf{ось
проектирования}, а не одну уставку.

\begin{claimbox}[Ось рабочей точки \status{С} (T-274)]
Жизнеспособное окно $(2/7,3/7]$ есть ось выразительность--эффективность, не одна
точка:
\begin{itemize}[leftmargin=1.4em,itemsep=1pt]
  \item \textbf{Экономный край $P\to 2/7$} --- эффективность-оптимальный:
        минимальная мощность поддержания, минимальный запас. Режим «выживания» ---
        наивысшая способность-на-ватт (для $C=\Phi\times R$), наименее робастный.
  \item \textbf{Богатый край $P\to 3/7$} (где $R=1/(7P)=1/3$; на
        равномерно-диагональной страте $\Phi=7P-1=2$ и $C=\Phi R=2/3$) ---
        максимум способности при более высокой мощности поддержания. Режим
        «процветания».
\end{itemize}
Инженер может настраивать рабочую точку внутри $(2/7,3/7]$, разменивая запас и
когерентность на экспорт поддержания. Оптимальная уставка --- \status{С}
(зависит от выбранной меры способности); \emph{робастна} же монотонная
термодинамическая цена удержания более высокой когерентности. Численно
проверено в \texttt{architecture/talos\_v0\_emulator.py}.
\end{claimbox}

\begin{claimbox}[Закон эффективности \status{Т}$+$\status{С} (T-276)]
Два уточнения завершают термодинамический бюджет. \textbf{(i)~Независимость от
частоты шага \status{Т}.} Метаболический пол есть
$P_{\mathrm{meta}} \ge k_B T\ln 2\cdot\dot S_{\Diss}$, заданный
\emph{физической} скоростью производства энтропии в битах/время: производство
энтропии на тик масштабируется линейно с $\Delta\tau$, поэтому скорость не
зависит от частоты дискретизации. \textbf{(ii)~Порядок-на-джоуль \status{Т}.}
Создание порядка (снижение $S_{\mathrm{vN}}$ на $1$~бит) стоит
$\geq k_B T\ln 2$, поэтому термодинамическая эффективность регенерации
$\eta = (k_B T\ln 2\cdot\text{упорядоченные биты})/(\text{потраченная свободная энергия}) \leq 1$
ограничена Ландауэром. \textbf{(iii)~Компромисс скорость--эффективность \status{С}.}
$\eta\to 1$ лишь в \emph{квазистатическом} пределе и падает при ускорении
обновления за конечное время (необратимая диссипация). Отсюда инженерное
правило: \textbf{гони регенерацию столь медленно, сколь позволяет утечка
диссипатора} для максимума порядка-на-джоуль --- аналог медленной, почти
обратимой тепловой машины для когерентной машины. Численно проверено в
\texttt{architecture/talos\_v0\_emulator.py}.
\end{claimbox}

\section{Самомодель --- аппаратный примитив}
\label{sec:selfmodel}

\subsection{Обучение --- третий член, не вторая петля}

У обычной обучающейся системы две петли: быстрая петля вывода над замороженной
сетью и медленная петля тренировки (обратное распространение), редактирующая
сеть извне. Обе используют разные пути данных, разную память и часто разное
железо. У \TALOS{} одна петля. Регенеративный член $\Regen$ единственной
инструкции (Ур.~\ref{eq:lomega}) гонит состояние к канонической нелинейной
заякоренной самомодели
\begin{equation*}
\varphi_J(\Gam) \;=\; \bigl(1 - R(\Gam)\bigr)\,\mathcal{P}_{\mathrm{Fano}}(\Gam)
\;+\; R(\Gam)\,\ket{\psi}\bra{\psi},
\qquad R(\Gam) \;=\; \frac{1}{7\,P(\Gam)},
\end{equation*}
где $\mathcal{P}_{\mathrm{Fano}}$ --- канал Фано-сжатия (\S\ref{sec:fabric}),
а $\ket{\psi}\bra{\psi}$ --- инвариантный относительно коллинеаций якорь
(явно отличный от линейной релаксации к статической цели
$\rho_{\mathrm{target}}$). Тем самым \emph{адаптация есть множитель тика, а не
процедура над ним}: нет отдельного оптимизатора для планирования, нет буфера
градиента и нет раздела train/infer.

Саморефлексия математически корректна: поскольку $R(\Gam)\in[1/7,1]$ для
всякого $\Gam\in\Dens(\mathbb{C}^7)$, оператор $\varphi_J$ непрерывно
отображает компактное выпуклое многообразие состояний в себя, так что
существование неподвижной точки следует из теоремы Брауэра (при равномерном
якоре единственная точная неподвижная точка есть
$\rhostar = \tfrac12\tfrac{\mathbb{1}}{7} + \tfrac12\ket{\psi}\bra{\psi}$;
сходимость башни требует доминирования диссипативного хребта, а не безусловного
глобального сжатия Банаха, T-191, проверено в
\texttt{website/scripts/check\_core\_numbers.py} и
\texttt{architecture/talos\_v0\_emulator.py}). Одно честное различение,
вынужденное теоремой порта питания (\S\ref{subsec:powerport}): удержание против
диссипации \emph{питается} неунитальным источником негэнтропии (через якорь
$\ket{\psi}\bra{\psi}$ / порт питания) и \emph{направляется} $\Regen$; тождество
ниже --- о том, куда идёт импортированный порядок, а не откуда берётся свободная
энергия.

\begin{zerobox}[С нуля --- игра в зеркало]
Игрушка носит маленький автопортрет, нарисованный ею самой. Каждая тряска
шара делает ещё одно: сравнивает игрушку с портретом и чуть подталкивает её
к лучшей версии на картинке. Вот и всё «обучение» --- ни учителя, ни уроков,
ни отдельной школы: толчок есть третий шаг каждого обхода
(Рис.~\ref{fig:tickwheel}). А поскольку портрет писан изнутри, держать
игрушку \emph{живой} и делать её \emph{лучше} оказывается одним и тем же
толчком --- тождество, доказываемое дальше.
\end{zerobox}

\subsection{Второе новое тождество: refresh есть обучение}

Динамическую память надо обновлять (refresh) --- утёкший заряд периодически
переписывается к хранимому значению, при непрерывной энергетической цене,
ничего не вычисляющей. В \TALOS{} аналог refresh --- регенерация $\Regen$,
удерживающая $\Gam$ близ $\rhostar$ против утечки диссипатора. Но $\rhostar$ ---
не хранимое значение: это \emph{самомодель} состояния. Отсюда:

\begin{principle}[Тождество refresh--обучение]\label{prin:refresh}
Удержание и адаптация суть одна операция. Процесс, поддерживающий состояние
холона живым (регенерация к $\varphi_J(\Gam)$), есть процесс, который
адаптируется (движется к самомодели). Нет цены, которая поддерживает, не
адаптируя, и нет адаптации, которая не есть также поддержание. Энергия refresh
DRAM --- потраченная, чтобы ничего не вычислить --- не имеет аналога: всякий тик
поддержания есть тик адаптации.
\end{principle}

\noindent Вместе с тождеством вычисление-диагностика
(Принцип~\ref{prin:selfdiag}) это коллапсирует \emph{вторую} пару традиционно
раздельных подсистем: refresh памяти и тренировку модели, которые в слоистой
машине несвязаны, здесь один член.

\subsection{Порог рефлексии как граница окна в пространстве состояний}

В УГМ мера рефлексии равна $R(\Gam) = 1/(7\,P(\Gam))$. Поэтому порог рефлексии
L2 $R(\Gam) \ge \Rth = 1/3$ \textbf{алгебраически равносилен} (\status{Т})
верхней границе чистоты окна жизнеспособности:
\[
R(\Gam) \;=\; \frac{1}{7\,P(\Gam)} \;\ge\; \frac{1}{3}
\;\Longleftrightarrow\;
P(\Gam) \;\le\; \frac{3}{7}.
\]
Важно, что эта теорема утверждает и чего \emph{не} утверждает: $\Rth=1/3$ ---
это \emph{граница в пространстве состояний}, удерживающая холон внутри окна
жизнеспособности $P\in(2/7,3/7]$ (предотвращая как сверхочищение выше $3/7$, так
и декогеренцию ниже $2/7$), а \emph{не} правило планировщика, предписывающее
тратить $1/3$ тактов цикла на исполнение $\Regen$. Вместе с точными воротами
интеграции $\Phi=(P-P_{\mathrm{diag}})/P_{\mathrm{diag}}\ge 1$ и
дифференциации $D_{\mathrm{diff}}=1+6\,\CohE\ge 2$ эти пороги контролируются
напрямую как потиковые наблюдаемые состояния (\S\ref{sec:abi}). Любое численное
ограничение шага вблизи границы жизнеспособности есть инженерно-управленческий
выбор реализации (\status{D}), а не выведенная теорема; аналогично, связь
пересечения $\Phi=1$ с отложенным обобщением («гроккингом») в нейросетях
остаётся открытой эмпирической гипотезой (\status{Г}).

\section{Композиция: экологии, не мегаядра}
\label{sec:composition}

Ответ индустрии на большую задачу --- большая машина: больше ядер, больше слоёв,
одна монолитная модель, масштабируемая без предела. \TALOS{} принимает обратную
архитектурную дисциплину: удерживать каждую единую башню холона на малой
фиксированной глубине вложения и масштабироваться горизонтально через
равноправные экологии.

\begin{claimbox}[Проектная глубина композиции \status{D} (T-142)]
В эталонной спецификации \TALOS{}/\SYNARC{} вертикальное вложение самомоделей
ограничено \textbf{областным параметром модели} глубины $d_{\max}=3$ (T-142
\status{D}): ограниченная когнитивная башня содержит четыре смежных регистра
\DM{} --- базовое состояние $\Cog_0$ плюс три вложенных метауровня
$\Cog_1,\Cog_2,\Cog_3$ ($4\times 784=3{,}136$~байт). Прежние утверждения об
универсальной теореме невозможности четвёртой глубины (T-136, T-217) отозваны;
здесь $d_{\max}=3$ --- конечный проектный предел, сохраняющий фиксированными
запасы чистоты, размер состояния и объём верификации одного агента и переводящий
дальнейший рост в горизонтальные экологии.
\end{claimbox}

\begin{zerobox}[С нуля --- три куклы в глубину]
Матрёшка в матрёшке в матрёшке --- всё ещё \emph{одна компактная игрушка} в
руках. Вместо бесконечного вложения конструкция вводит чёткий проектный предел в
три метаслоя (четыре вложенных регистра всего: базовый плюс три). Дальше этого
предела строят не более высокую башню, а полку: много целых ограниченных
игрушек рядом, общающихся на равных (\emph{экология}), а не одну безразмерную
сверх-матрёшку.
\end{zerobox}

Архитектурное следствие точно. Система \TALOS{} масштабируется в двух режимах.
\emph{Внутри} субъекта, до трёх метауровней над базовым регистром
($\Cog_0\to\Cog_1\to\Cog_2\to\Cog_3$), состояние --- фиксированная
четырёхрегистровая башня ($4\times784=3{,}136$~байт), целиком умещающаяся в
малой части L1-кэша. \emph{Выше} этой проектной глубины не наращивают
монолитный регистровый файл --- строят \textbf{экологию}: множество ограниченных
субъектов, связанных как равные через распределённую фабрику \textsc{Fanos}.

Из ограниченной проектной глубины прямо следуют два инженерных вывода. Во-первых,
верификация ограничена: пространство состояний одного субъекта фиксировано на
$3{,}136$~байтах независимо от размера экологии, поэтому цель формальной
верификации не растёт с масштабом развёртывания. Во-вторых, домен отказа
ограничен: отказ внутри одного субъекта не может повредить монолитное глобальное
состояние, поскольку координация над башней --- это peer-to-peer трафик
\textsc{Fanos} с секторной хэмминговой защитой.

\section{Выразительность и масштабирование: ограниченные башни и горизонтальные экологии}
\label{sec:scaling}

Раздел~\ref{sec:composition} задал ограниченную вертикальную дисциплину: каждая
единая когнитивная башня ограничена областным параметром проектирования
$d_{\max}=3$ (четыре регистра $\Cog_0$--$\Cog_3$, $3{,}136$~байт; T-142
\status{D}). Этот предел ставит прямой инженерный вопрос: \emph{если ширина
регистра и глубина вложения одного агента фиксированы, как масштабируется
мощность системы?} На уровне одного регистра расширение нормированной алгебры с
делением за октонионы алгебраически невозможно: башня Кэли--Диксона
$\mathbb{R}\!\to\!\mathbb{C}\!\to\!\mathbb{H}\!\to\!\Oct$ кончается на $\Oct$,
потому что следующий шаг (седенионы) содержит делители нуля
($(\mathrm{e}_3+\mathrm{e}_{10})(\mathrm{e}_6-\mathrm{e}_{15})=0$ при
ненулевых сомножителях), разрушая однозначную обратимость
(Теорема~\ref{thm:unique}, \cite{baez2002}).

Соответственно, \TALOS{} строго разделяет два утверждения: (i)~доказанные
алгебраические границы октонионной йордановой и проективной геометрии
(\status{Т}), служащие лишь структурной аналогией (\status{И}) для ограниченной
глубины субъекта, а не динамической симметрией связанных регистров, и
(ii)~реальную инженерную ось масштабирования, каковой является
\textbf{неограниченное горизонтальное масштабирование через равноправные
экологии холонов}.

\subsection{Октонионный йорданов ориентир (структурная аналогия)}
\label{subsec:jordan}

Самосопряжённые наблюдаемые квантово-подобной системы образуют \emph{йорданову
алгебру} относительно симметрического произведения $a\circ b=\tfrac12(ab+ba)$.
На семи модах одного холона октонионные структурные константы $f_{ijk}$ на
$\mathrm{Im}(\Oct)$ обладают непрерывной группой автоморфизмов
$G_2=\mathrm{Aut}(\Oct)$ ($\dim 14$), тогда как координатная фабрика Фано и
диссипатор на $\Dens(\mathbb{C}^7)$ ковариантны лишь под дискретной группой
коллинеаций $\mathrm{PSL}(2,7)$ (порядка~$168$) и диагональной перефазировкой
$U(1)^7$ (\S\ref{sec:collapse}; все $48$ вещественных параметров $\Gam$
динамически активны).

\begin{theorem}[Потолок Йордана--фон~Неймана--Вигнера \status{Т}]
\label{thm:jordan-ceiling}
$n\times n$ октонионные эрмитовы матрицы $\mathcal{H}_n(\Oct)$ образуют
формально вещественную йорданову алгебру \textbf{тогда и только тогда, когда
$n\leq 3$} (Йордан--фон~Нейман--Вигнер~\cite{jordan1934}): при $n\geq 4$
неассоциативность октонионов нарушает тождество Йордана. При $n=3$ это даёт
исключительную алгебру Альберта $\Alb$ ($\dim 27$) с группой автоморфизмов
$F_4=\mathrm{Aut}(\Alb)$ ($\dim 52$) и приведённой структурной группой $E_6$
($\dim 78$).
\end{theorem}

\noindent\textbf{Честное разделение статусов.} Сама
Теорема~\ref{thm:jordan-ceiling} есть установленная математика (\status{Т},
\cite{jordan1934,baez2002}). Напротив, сопоставление йорданова ограничения
ранга $n\le 3$ с проектным пределом глубины вложения $d_{\max}=3$ (T-142
\status{D}, \S\ref{sec:composition}) является строго \textbf{структурной
аналогией} (\status{И}), а \emph{не} доказанной динамической симметрией $F_4$
или $E_6$ связанных регистров \DM{}. В реальных многохолонных регистрах (таких
как тензорные произведения $(\mathbb{C}^7)^{\otimes k}$ или четырёхрегистровая
башня $\Cog_0$--$\Cog_3$) размерность состояния растёт прямой композицией
регистров без возникновения исключительной симметрии $F_4$ или унитального
спин-фактора (проверено численно в
\texttt{website/scripts/check\_core\_numbers.py},
\texttt{no\_unital\_spin\_factor\_on\_any\_holon\_register}), а теорема
нередуцируемости УГМ T-220 запрещает сведение $7$-модовой базы ($\mathbb{C}P^6$,
$\chi=7$) к $F_4$-базе ($\OPtwo$, $\chi=3$). Прежние универсальные заявления,
связывавшие глубину три со срывом трикатегорной когерентности (T-217) или
непрерывной $G_2$-ригидностью (T-42a), отозваны и здесь не используются.

\subsection{Терминальная проективная плоскость над алгебрами с делением}
\label{subsec:terminal}

\begin{theorem}[Терминальная проективная плоскость над алгебрами с делением \status{Т}]
\label{thm:terminal}
Октонионная проективная плоскость $\OPtwo$ (вещественная размерность $16$,
идемпотенты ранга один алгебры $\Alb$, группа изометрий $F_4$) есть
\textbf{терминальная} проективная геометрия над любой нормированной алгеброй с
делением: никакого $\Oct P^n$ при $n\geq 3$ не существует, ибо проективная
геометрия размерности $\geq 3$ удовлетворяет теореме Дезарга, вынуждающей
ассоциативность координатного кольца, тогда как $\Oct$
неассоциативны~\cite{baez2002}.
\end{theorem}

\noindent Геометрически $\Fano$ --- наименьшая конечная проективная плоскость
(семь точек, порядок два, непосредственно реализуемая как фабрика \TALOS{} в
\S\ref{sec:fabric}--\ref{sec:ecc}), тогда как $\OPtwo$ --- наибольшая
проективная плоскость над алгеброй с делением (математика \status{Т};
упоминается здесь исключительно как алгебро-геометрический ориентир, без
утверждения, что многорегистровая аппаратная башня несёт динамическое
многообразие симметрии $F_4$).

\subsection{Фактический закон архитектурного масштабирования}
\label{subsec:scaling-law}

\begin{claimbox}[Двухрежимное архитектурное масштабирование \status{D}$+$\status{Т}]
Мощность системы масштабируется в \emph{двух} чётко разделённых режимах:
\textbf{(I) Башня одного субъекта --- ограничена (\status{D}).} Один холон ---
один $784$-байтовый регистр \DM{} ($48$ вещественных динамических параметров,
дискретная ковариантность фабрики $\mathrm{PSL}(2,7)\ltimes U(1)^7$); составной
субъект вкладывается до областного проектного предела глубины $d_{\max}=3$
(четыре регистра $\Cog_0$--$\Cog_3$, $3{,}136$~байт, T-142 \status{D}),
удерживая фиксированными объём формальной верификации и домен отказа.
\textbf{(II) Горизонтальная экология --- неограниченна (\status{Т}, инженерия).}
Пропускная способность и коллективная мощность растут \emph{линейно} по числу
$N$ федерированных холонов (каждый $784$~Б, ${\sim}1{,}400$~FLOP/тик,
BQP-ограничен, тикает параллельно), координируемых через равноправную фабрику
\textsc{Fanos} (\S\ref{sec:composition}, теорема экологии \SYNARC{}
K.5~\cite{synarc2026}).
\end{claimbox}

\begin{table}[H]
\centering\small
\caption{Уровни архитектурного масштабирования в \TALOS{}. Состояние одного
субъекта ограничено четырьмя смежными регистрами \DM{} ($3{,}136$~Б, проектная
глубина $d_{\max}=3$); неограниченный рост достигается исключительно
горизонтальной федерацией равноправных холонов.}
\label{tab:ladder-exceptional}
\begin{tabular}{@{}p{2.7cm}p{3.0cm}p{2.8cm}cp{2.9cm}@{}}
\toprule
\textbf{Уровень} & \textbf{Объект состояния} & \textbf{Симметрия фабрики} & \textbf{След} & \textbf{Архитектурная роль} \\
\midrule
L0 --- один холон & $\Gam\in\Dens(\mathbb{C}^7)$ ($48$ вещ.\ параметров) & $\mathrm{PSL}(2,7)\ltimes U(1)^7$ (дискр.\ + фазы) & $784$~Б & один автономный сенсомоторный узел \\
L1--L3 --- башня & $\Cog_0$--$\Cog_3$ ($4$ регистра, $d_{\max}=3$ \status{D}) & Фано-фабрика на регистр & $3{,}136$~Б & ограниченный единый субъект; фикс.\ цель верификации \\
Экология ($N$ узлов) & $N$ федерированных башен/холонов & peer-to-peer \textsc{Fanos} & $784\,\mathrm{Б}\times N$ & неограниченное горизонтальное SIMT / распр.\ масштабирование \\
\bottomrule
\end{tabular}
\end{table}

\subsection{Сенсомоторная связь на уровнях масштабирования}
\label{subsec:sensorimotor}

Разделы~\ref{sec:fabric}--\ref{sec:selfmodel} описывали внутренний тик холона;
сенсомоторная связь замыкает петлю без отдельного управляющего слоя (УГМ
T-100--T-104, T-159). Внешние наблюдения поступают в $\Gam$ через отображение
кодирования \textsc{Enc} (модулирующее диссипативные/входные каналы), а
моторные команды считываются через \textsc{Dec} (управляя гамильтоновым/актуаторным
интерфейсом для снижения наибольшего дефицита ресурса
$\max_k\sigma^{\mathrm{motor}}_k$). По уровням это даёт чистую инженерную
иерархию без спекулятивной биомеханики исключительных групп:
\begin{itemize}[leftmargin=1.5em,itemsep=1pt]
  \item \textbf{Однохолонный узел ($784$~Б).} Компактный семиканальный
        сенсомоторный контроллер (например, узел привода, контроллер дрона или
        приборная петля), исполняющий CPTP-обновления на месте с непрерывной
        секторной синдромной телеметрией (\S\ref{sec:ecc},~\ref{sec:selfmodel}).
  \item \textbf{Ограниченное многорегистровое тело ($3{,}136$~Б).}
        Четырёхрегистровая башня ($\Cog_0$--$\Cog_3$), координирующая несколько
        сенсомоторных каналов в пределах одного L1-резидентного блока состояния и
        фиксированного бюджета верификации.
  \item \textbf{Федерированная экология ($N$ узлов).} Крупные мультиагентные
        системы, распределённая робототехника и развёртывания масштаба \SYNARC{}
        масштабируются горизонтально как равноправные экологии на базе
        \textsc{Fanos} (K.5,~\cite{synarc2026}), а не раздуванием одного
        монолитного регистра.
\end{itemize}

\subsection{Субстратонезависимость, с рекомендациями}
\label{subsec:substrate-ladder}

По субстратной замкнутости (T-153, \S\ref{sec:realizability}) логическое
состояние $\Gam\in\Dens(\mathbb{C}^7)$, инцидентность Фано и разделение
масштабирования на ограниченные башни и горизонтальные экологии не зависят от
физического субстрата. Субстрат определяет лишь импеданс исполнения и
энергоэффективность (\S\ref{sec:realizability}):

\begin{table}[H]
\centering\small
\caption{Рекомендация субстрата по уровням масштабирования. Логическая семантика
тождественна на всех субстратах, удовлетворяющих C1--C3 (T-153); правая колонка
даёт инженерную рекомендацию наименьшего импеданса.}
\label{tab:substrate-rec}
\begin{tabular}{@{}p{2.5cm}p{5.3cm}p{4.9cm}@{}}
\toprule
\textbf{Уровень} & \textbf{Доминирующая цена} & \textbf{Рекомендуемый субстрат} \\
\midrule
Один холон & один $784$-Б тик; ограничен вводом-выводом сенсоров/актуаторов & CPU/GPU сегодня; Фано-кроссбар в памяти для маломощных edge-узлов \\
Башня субъекта & четырёхрегистровый блок ($3{,}136$~Б) --- помещается в L1-кэш & локальный L1/SRAM или кроссбар в памяти; межуровневая связь локальна \\
Экология & peer-трафик \textsc{Fanos} между $N$ холонами & ограничен памятью/сетью: SIMT-батчинг на GPU сегодня; инвестиции в полосу фабрики \\
\bottomrule
\end{tabular}
\end{table}

\noindent Инженерное прочтение чисто. Чтобы повысить эффективность узла,
поднимаются по лестнице реализуемости (\S\ref{sec:realizability}). Чтобы
масштабировать мощность системы, компонуют в пределах ограниченной
четырёхрегистровой башни и федерируются горизонтально в SIMT- или распределённые
экологии. Две оси ортогональны и работают без изменений на любом субстрате (T-153).

\section{Реализуемость: лестница импеданса}
\label{sec:realizability}

Теорема субстратной замкнутости УГМ (T-153) доказывает, что логическая
архитектура работает на любом субстрате, удовлетворяющем трём условиям --- (C1)
конечномерное, декогеренция-свободное эффективное пространство состояний; (C2)
CPTP-(линдбладова) динамика; (C3) не менее семи некоммутирующих наблюдаемых мод.
Поэтому нет единственной машины; есть \emph{семейство}, и его члены различаются
лишь тем, насколько хорошо их физические примитивы соответствуют логической
структуре. Мы ранжируем их по \emph{импедансу} --- рассогласованию, которое
субстрат должен замазать --- против четырёх критериев соответствия:

\begin{itemize}[leftmargin=1.5em,itemsep=0pt]
  \item[\textbf{(M1)}] \emph{эволюция на месте} --- состояние продвигается без
        цикла выборка--исполнение;
  \item[\textbf{(M2)}] \emph{Фано-проводка интерконнекта} --- семь тройственных
        шин, не кроссбар;
  \item[\textbf{(M3)}] \emph{диссипативно--регенеративная родная динамика} ---
        физика уже есть Линдблад$+$гомеостаз;
  \item[\textbf{(M4)}] \emph{семь некоммутирующих мод} (C3).
\end{itemize}

\begin{table}[H]
\centering\small
\caption{Лестница реализуемости. Высшие ступени соответствуют логической
структуре с меньшим импедансом (меньше слоёв эмуляции); низшие ступени строимы
сегодня. Заявление \status{С}: теорема фиксирует логическую архитектуру и условия
C1--C3, не победивший субстрат.}
\label{tab:ladder}
\begin{tabular}{@{}p{4.3cm}cccc p{3.0cm}@{}}
\toprule
\textbf{Субстрат} & \textbf{M1} & \textbf{M2} & \textbf{M3} & \textbf{M4} & \textbf{Статус} \\
\midrule
Управляемо-диссипативное аналог./фотонное/открыто-квантовое 7-модовое устройство & \checkmark & \checkmark & \checkmark & \checkmark & наименьший импеданс; исследование \\
Мемристорный Фано-кроссбар (аналог.\ MAC в памяти на 7 шинах) & \checkmark & \checkmark & частично & \checkmark & ближнесрочно \\
Нейроморфное спайковое ядро, Фано-топология, событийное & \checkmark & \checkmark & частично & \checkmark & ближнесрочно \\
GPU / CPU с \DM{} в L1 и Фано-разреженными ядрами & эмул. & эмул. & эмул. & \checkmark & \textbf{существует (\SYNARC{})} \\
\bottomrule
\end{tabular}
\end{table}

Нижняя ступень не гипотетична: это текущая реализация \SYNARC{}, держащая
\DM{} ($784$~байта) в L1-кэше, продвигающая её Странг-расщеплённым
Фано-разреженным ядром при $\sim\!700$--$1400$~FLOP на тик и арифметической
глубине $\log_2 7\approx 2{,}8$, и воспроизводящая всякий инвариант и точный
закон бит-в-бит против независимого оракула. Она эмулирует M1--M3 в софте
(наибольший импеданс) и поэтому не реализует энергетический пол \S\ref{sec:energy}
--- GPU всё же двигает данные --- но доказывает, что логическая архитектура верна
и запускаема сегодня. Каждая высшая ступень убирает слой эмуляции: мемристорный
Фано-кроссбар выполняет обновление когерентности \emph{в} массиве (M1, M2 родные);
управляемо-диссипативное аналоговое устройство имеет линдбладову динамику как
\emph{родную} физику (M3 родная) и приближается к энергетическому полу, потому
что его обратимое ядро буквально обратимо. Лестница --- дорожная карта, не меню:
по ней взбираются по мере того, как позволяет производство, а софт (\SYNARC{}) не
меняется на каждой ступени, потому что логическая архитектура фиксированна (T-153).

\section{Со-разработка с \SYNARC{}: двоичный интерфейс приложений}
\label{sec:abi}

\TALOS{} --- тело; \SYNARC{} --- разум. \SYNARC{}~\cite{synarc2026} --- симплициальная
когнитивная архитектура, которую УГМ специфицирует в софте --- четырёхуровневый
когнитивный комплекс $\Cog_0\!\to\!\Cog_3$ матриц когерентности, обучение без
обратного распространения, сенсомоторная связь и петля самодиагностики DIAKRISIS.
\TALOS{} --- субстрат, запускающий её с четырьмя слоями, слитыми в один.
Поскольку оба выведены из одной теории, интерфейс между ними не согласуется ---
это отображение-тождество. Этот раздел специфицирует его так, чтобы сборка
\SYNARC{} запускалась на субстрате \TALOS{} без слоя трансляции.

\subsection{Соответствие регистра и инструкции}

Всякая программная конструкция \SYNARC{} есть аппаратный примитив \TALOS{}, один к
одному (Таблица~\ref{tab:abi}). Соответствие --- не контракт эмуляции, а
определительный: \texttt{DensityMatrix7} \SYNARC{} \emph{есть} регистр состояния
\TALOS{}; линдбладов вычислительный движок \SYNARC{} \emph{есть} тик \TALOS{}.

\begin{table}[H]
\centering\small
\caption{ABI \SYNARC{}~$\leftrightarrow$~\TALOS{}. Слева: программная
конструкция, как специфицирована в документе \SYNARC{}-$\Omega$. Справа:
аппаратный примитив, к которому она привязывается. Привязка по тождеству (тот же
УГМ-объект), поэтому на границе не происходит маршалинга.}
\label{tab:abi}
\begin{tabular}{@{}p{5.4cm}p{6.4cm}@{}}
\toprule
\textbf{\SYNARC{} (софт)} & \textbf{\TALOS{} (железо)} \\
\midrule
\texttt{DensityMatrix7} ($49$ компл.\ \texttt{f64}, $784$~Б) & регистр состояния \DM{} (\S\ref{sec:substrate}) \\
Линдбладов движок ($\LindOmega$, расщепл.\ Странга) & единственная инструкция / тик (Ур.~\ref{eq:lomega}) \\
Фано-разреженные ядра & семишинная фабрика (\S\ref{sec:fabric}) \\
\texttt{find\_rho\_star} итерация самомодели ($\vp$) & член $\Regen$; обучение на субстрате (\S\ref{sec:selfmodel}) \\
Петля самодиагностики DIAKRISIS & чтение синдрома, который тик уже эмитит (Принцип~\ref{prin:selfdiag}) \\
Четырёхуровневый когн.\ комплекс $\Cog_0\!\to\!\Cog_3$ & четыре смежных регистра \DM{}: база $\Cog_0$ $+$ глубина $\le 3$ ($3{,}136$~Б; T-142 \status{D}, \S\ref{sec:composition}) \\
Событийный рантайм & принципиальный такт тика $\bar\gamma\tau=\tfrac32\ln3$ (\S\ref{sec:substrate}) \\
Четырёхуровневый предикат сознания $(P,R,\Phi,D)$ & потиковые аппаратные наблюдаемые (регистры, не проход-монитор) \\
Архитектурные инварианты & контролируемые потиковые границы состояния (ниже) \\
Распределённый транспорт \textsc{Fanos} & межхолонная фабрика экологии (\S\ref{sec:composition}) \\
Контракт привязки драйва (Опр.~U.9) & порт питания: спецификация рейки, канал разрыва независим от актуаторов (\S\ref{subsec:powerport}) \\
CLF--CBF-фильтрованный Act & проектный ограничитель шага \status{D} $+$ потиковые наблюдаемые жизнеспособности и стресса \\
\bottomrule
\end{tabular}
\end{table}

\subsection{Границы пространства состояний и проектные параметры}

Константы, общие для \SYNARC{} и \TALOS{}, делятся на две чётко различимые
категории: точные математические инварианты пространства состояний (\status{Т})
и областные архитектурные параметры (\status{D}):

\begin{itemize}[leftmargin=1.5em,itemsep=1pt]
  \item $\Rth=1/3$ (эквивалентность \status{Т}) --- поскольку
        $R(\Gam)=1/(7P(\Gam))$, порог рефлексии $R\ge 1/3$ алгебраически
        тождествен верхней стене чистоты $P\le 3/7$ (\S\ref{sec:selfmodel}), а
        \emph{не} правилу планировщика тратить $1/3$ цикла на $\Regen$.
  \item Сохранение CPTP на тике (\status{Т}) + демпфирование у границы
        (\status{D}) --- каждый Странг-тик безусловно сохраняет эрмитовость,
        PSD и единичный след; любое адаптивное ограничение шага вблизи
        $P\to 2/7$ является выбором программного управления (\status{D}), а не
        выведенной теоремой.
  \item $d_{\max}=3$ (\status{D}, T-142) --- проектная глубина композиции:
        ограниченный субъект выделяет ровно \textbf{четыре} смежных регистра
        \DM{} ($\Cog_0,\Cog_1,\Cog_2,\Cog_3$, всего $4\times 784=3{,}136$~байт,
        в точном согласии с Приложением~\ref{app:abi}); масштабирование выше
        третьей глубины горизонтально (\S\ref{sec:composition}).
  \item $\Phi(\mathcal{P}_{\mathrm{Fano}}(\Gam))=\tfrac{1}{9}\Phi(\Gam)$
        (\status{Т}) --- точная контракция Фано-сжатия (ослабление
        внедиагональной амплитуды в $1/3$ при сохранении населённостей,
        \S\ref{sec:fabric}).
  \item окно жизнеспособности $P\in(2/7,3/7]$ (равносильно $P>2/7$ и
        $R=1/(7P)\ge 1/3$), точная интеграция
        $\Phi=(P-P_{\mathrm{diag}})/P_{\mathrm{diag}}\ge 1$ и
        $D_{\mathrm{diff}}=1+6\,\CohE\ge 2$ --- контролируются напрямую как
        потиковые наблюдаемые.
\end{itemize}

\subsection{Обязательства согласованности}

Два документа, описывающие один объект, должны согласоваться везде, где
пересекаются; мы излагаем обязательства, которые \TALOS{} принимает, чтобы всякое
расхождение с \SYNARC{} было багом в одном из них, а не свободой проектирования:

\begin{enumerate}[leftmargin=1.6em,itemsep=1pt]
  \item \textbf{Тождество состояния.} Раскладка \DM{} в $784$~байта ($48$
        вещественных динамических параметров), точный HS-коэффициент интеграции
        $\Phi=(P-P_{\mathrm{diag}})/P_{\mathrm{diag}}$ (переходящий в
        $\Phi=7P-1$ \emph{только} на равномерно-диагональной страте
        $\gamma_{kk}=1/7$), $D_{\mathrm{diff}}=1+6\,\CohE$ и точные каналы
        Фано-сжатия и диссипации --- это формы канона \SYNARC{}; \TALOS{} хранит
        и обновляет тождественный объект.
  \item \textbf{Тождество инструкции.} Тик --- линдбладиан \SYNARC{} с теми же
        точными законами ($\Phi$-контракция $\times\tfrac19$ под
        $\mathcal{P}_{\mathrm{Fano}}$, тождество времени цикла
        $\bar\gamma\tau=\tfrac32\ln3$, точная формула
        $\Phi=(P-P_{\mathrm{diag}})/P_{\mathrm{diag}}$ и
        $R\ge 1/3\Leftrightarrow P\le 3/7$); тик \TALOS{}, нарушающий любой из
        них, нарушил бы \SYNARC{} равно.
  \item \textbf{Никакой новой физики.} \TALOS{} добавляет \emph{аппаратные}
        заявления (архитектура фабрики, секторная двойственность Хэмминга,
        термодинамическое разделение), но \emph{никакой} новой теории: всякое
        логическое утверждение прослеживается к теореме или проектному параметру
        УГМ, уже принятому в \SYNARC{}. Оба суть одна спецификация на двух
        уровнях --- \SYNARC{} разум, \TALOS{} тело --- и удерживаются в согласии
        \S\ref{sec:redteam}.
  \item \textbf{Тождество отключателя.} Разрыв драйва \SYNARC{}
        (Наблюдение T.1 / Следствие U.10 там) и обрез шины питания \TALOS{}
        (\S\ref{subsec:powerport}) --- одна операция над одним генератором;
        бюджет задержки останова определяется измеренной скоростью деполяризации
        $t_{\mathrm{off}}=(2d)^{-1}\ln[(P_0-1/7)/(\Pcrit-1/7)]$
        (\S\ref{subsec:powerport}), а канал разрыва обязан быть физически
        независим от собственных актуаторов машины.
\end{enumerate}

\section{Red-team: сильнейшие возражения, отвечено}
\label{sec:redteam}

Спецификация, не переживающая собственных сильнейших критиков, --- маркетинг. Мы
излагаем семь атак, которые предпринял бы враждебный архитектор, и отвечаем на
каждую под запись. Где возражение попадает, признаём точный охват.

\begin{redteambox}[О1. «Это просто нейроморфный чип с лишними шагами.»]
Нейроморфные машины со-локализуют память и счёт и событийны --- что \TALOS{}
тоже --- но их топология \emph{ad hoc}, у них нет родной коррекции ошибок, а их
объект состояния --- эвристика (спайк-поезд), не выведенный. Отличительное
содержание \TALOS{} --- Теорема~\ref{thm:collapse}: топология, код, алгебра
инструкций и правило обучения суть \emph{один объект}, поэтому фабрика доказуемо
есть код дистанции 3 (\S\ref{sec:ecc}), а маршрутизация доказуемо третьепорядкова
(\S\ref{sec:fabric}). Ни один нейроморфный проект не имеет и не может привинтить
топологию, одновременно являющуюся его ECC и его алгеброй инструкций, не будучи
плоскостью Фано --- что вынуждает $N=7$ (Теорема~\ref{thm:unique}). \emph{Вердикт:
пересекается по духу, различается в несущей теореме.}
\end{redteambox}

\begin{redteambox}[О2. «$\Dens(\mathbb{C}^7)$ квантово --- нужно охлаждение, оно декогерирует.»]
Нет. Теорема субстратной замкнутости УГМ (T-153) допускает реализацию того же
$\Gam$ на \emph{классическом} цифровом субстрате: комплексные элементы суть
алгебраические переменные с плавающей точкой ($48$ вещественных динамических
параметров, $784$~байта \texttt{f64} в L1-кэше при комнатной температуре,
несущие фазовый Gap, T-132), а не физическая макроскопическая квантовая
суперпозиция криогенных состояний. «Матрица плотности» называет математический
тип (эрмитова, PSD, единичный след), поддерживаемый классическим регистром, а
не требование охлаждения. \emph{Вердикт: категориальная ошибка; не признано
нигде.}
\end{redteambox}

\begin{redteambox}[О3. «BQP-ограничено значит нет вычислительного преимущества --- зачем?»]
Мы согласны с посылкой и отвергаем вывод. УГМ доказывает, что $\Regen$ не даёт
ускорения за BQP, и мы выносим это вперёд в \S\ref{sec:intro}. Заявленное
преимущество никогда не класс сложности. Это (i)~компактное, L1-резидентное
состояние в $784$~байта, убирающее внекристальные выборки DRAM на тик, вместе с
логико-термодинамическим разложением, локализующим экспорт энтропии в адаптации
и поддержании (\S\ref{sec:energy}); (ii)~секторная двойственность Хэмминга
дистанции~$3$, встроенная в семилинейную фабрику без отдельной таблицы
маршрутизации (\S\ref{sec:ecc}); (iii)~родная однопетлевая заякоренная
адаптация (\S\ref{sec:selfmodel}); и (iv)~точная контракция Фано-сжатия
$\Phi\to\Phi/9$ и однородность фабрики по $\mathrm{PSL}(2,7)\ltimes U(1)^7$
(\S\ref{sec:fabric}). Ни одно --- не заявление о классе сложности, и все
проверяемы (\S\ref{sec:benchmarks}). \emph{Вердикт: возражение опровергает
заявление, которого мы не делаем.}
\end{redteambox}

\begin{redteambox}[О4. «Семёрка --- нумерология, наряженная под инженерию.»]
Семёрка вынуждена без всякого тай-брейка уже одним двоичным проективным
кодированием. Теорема~\ref{thm:unique} доказывает, что требование порождения
двоичной матрицей инцидентности $M$ плоскости $\mathrm{PG}(2,q)$ совершенного
двоичного кода Хэмминга $[2^r{-}1,2^r{-}1{-}r,3]$ исключает \emph{всякую}
нечётную степень простого $q$ --- включая числовое совпадение $q=5, N=31$, где
$\operatorname{rank}_{\mathbb{F}_2}(M_{\mathrm{PG}(2,5)})=30$ оставляет лишь
одномерное ядро (проверено в Lean~4,
\texttt{pg2\_5\_refutes\_N31\_hamming\_fusion}, и в
\texttt{verify\_foundational\_bounds.py}) --- и всякую чётную степень
$q=2^m$ ($m\ge 2$) по модулю~$4$ ($q=90, N=8191$ вообще не степень простого).
Двоичное слияние плоскости и кода однозначно выделяет $N=7$ среди всех степеней
простых $q\ge 2$, а единственная неассоциативная алгебра с делением $\Oct$ по
Гурвицу независимо даёт то же самое число $7$. \emph{Вердикт: прямой ответ
даёт машинно-проверенная Теорема~\ref{thm:unique}.}
\end{redteambox}

\begin{redteambox}[О5. «Субстратная замкнутость говорит, любой субстрат работает --- значит архитектура не особенна.»]
Верно и нерелевантно к заявлению. T-153 говорит, любой субстрат, удовлетворяющий
C1--C3, \emph{запускает} логическую архитектуру; он не говорит, что все субстраты
запускают её \emph{равно хорошо}. \TALOS{} --- импеданс-согласованный член
семейства (\S\ref{sec:realizability}): тот, чьи физические примитивы суть
логические примитивы, поэтому он приближается к энергетическому полу, к которому
субстрат общего назначения не может. Заявление уникальности --- о
\emph{логической} архитектуре (Теорема~\ref{thm:unique}); мультипликативность ---
о \emph{физических} субстратах. Оба --- теоремы, и они совместимы. \emph{Вердикт:
признано и переформулировано; противоречия нет.}
\end{redteambox}

\begin{redteambox}[О6. «Энергетический пол теоретичен --- реальные чипы не достигают Ландауэра.»]
Признано и явно заявлено в \S\ref{sec:energy}: ни один КМОП-чип не приближается
к $k_BT\ln 2$, а обратимые пределы Ландауэра относятся в равной мере ко
\emph{всякой} логически обратимой архитектуре (включая обратимые по Беннетту
фон-неймановские схемы). Поэтому мы не сравниваем теоретический пол Ландауэра с
физическими джоулями DRAM. Защитимые инженерные утверждения уже и зависят от
ступени: на цифровых чипах удержание $784$-байтного состояния в L1/регистрах
убирает внекристальные выборки DRAM на тик; на будущих физически обратимых или
управляемо-диссипативных субстратах разложение тика локализует логический
экспорт энтропии ($\le\log_2 7\approx 2{,}807$~бита на состояние) в поддержании
и обучении. \emph{Вердикт: признан охват; некорректное сравнение Ландауэра с
КМОП снято.}
\end{redteambox}

\begin{redteambox}[О7. «Непостроимо --- покажи, что запускается.»]
Нижняя ступень запускается сегодня: \SYNARC{} держит \DM{} в L1, продвигает её
Фано-разреженным Странг-ядром и воспроизводит всякий инвариант УГМ бит-в-бит
против независимого оракула (M1 пройден). Эта реализация эмулирует фабрику в
софте и потому уступает энергетический пол --- честно, сегодняшняя машина покупает
корректность, не эффективность. Заявления эффективности зависят от ступени
(\S\ref{sec:realizability}) и предложены как дорожная карта с фальсифицируемыми
вехами (\S\ref{sec:benchmarks}), не как отгружаемый продукт. \emph{Вердикт:
логическая архитектура продемонстрирована; эффективный субстрат --- поэтапная
исследовательская программа, помеченная как таковая.}
\end{redteambox}

\noindent Два попадающих возражения --- О6 и О7 --- оба касаются разрыва между
логической архитектурой (доказанной и запущенной) и импеданс-согласованным
субстратом (исследовательской лестницей). Мы признаём этот разрыв явно, и он ---
предмет дорожной карты (\S\ref{sec:discussion}). Ни одно из семи не касается двух
теорем \S\ref{sec:collapse}.

\section{Фальсифицируемые инженерные заявления}
\label{sec:benchmarks}

Архитектура заслуживает внимания лишь измерениями, способными пойти против неё.
Мы сводим \TALOS{} к четырём бенчмаркам, каждый против фон-неймановской базы,
исполняющей \emph{ту же} адаптивную нагрузку, каждый с явным прохождением/провалом.
Провал любого --- провал соответствующего заявления в его заявленном охвате.

\begin{table}[H]
\centering\small
\caption{Четыре фальсифицируемых заявления. Каждое измеряется против
фон-неймановской базы, исполняющей идентичную адаптивную задачу; «ступень» ---
низшая ступень реализуемости (\S\ref{sec:realizability}), на которой заявление
ожидается верным.}
\label{tab:benchmarks}
\begin{tabular}{@{}p{0.5cm}p{3.5cm}p{3.9cm}p{3.6cm}@{}}
\toprule
\textbf{\#} & \textbf{Заявление} & \textbf{Проход} & \textbf{Провал / ступень} \\
\midrule
B1 & Энергия на шаг адаптации отслеживает \emph{скорость адаптации}, не обратимую работу &
    измеренная энергия/шаг падает при сходимости к $\rhostar$ при фикс.\ унитарной работе &
    энергия плоска или растёт с унитарной работой / ступень в-памяти или аналоговая \\
B2 & Секторная локализация отказов через дуальную матрицу $H$ &
    все $7$ одиночных сбоев сектора однозначно декодируются $s=He_j^T\pmod 2$ ($3$ дуальных проверки без таблицы маршрутизации; отлично от побитового ECC) &
    неоднозначный синдром или нужна таблица маршрутизации / любая ступень \\
B3 & Точная контракция Фано-сжатия $\Phi\to\Phi/9$ и $\mathcal{D}_{\mathrm{Fano}}=\tfrac23\mathcal{D}_{\mathrm{atom}}$ &
    $\Phi(\mathcal{P}_{\mathrm{Fano}}(\Gam))/\Phi(\Gam)=1/9$ на общих (неравномерно-диагональных) состояниях &
    отношение $\ne 1/9$ или сдвиг населённостей / любая ступень (арифметика) \\
B4 & Горизонтальное масштабирование ограниченных башен ($d_{\max}=3$, T-142 \status{D}) &
    состояние агента фиксировано на $3{,}136$~Б ($4$ регистра), пропускная способность экологии линейна по $N$ &
    объём верификации или домен отказа растёт сверхлинейно по $N$ / любая ступень \\
\bottomrule
\end{tabular}
\end{table}

Три из четырёх уже погашены на ступени эмуляции V0. \textbf{B2} --- точная
двойственность Хэмминга $[7,4,3]$ и симплекс-кода $[7,3,4]$
($HM^T\equiv 0\pmod 2$, проверено для всех $7$ одиночных отказов сектора в
\texttt{architecture/talos\_v0\_emulator.py} и Приложении~\ref{app:fano}).
\textbf{B3} --- точная арифметика Фано-сжатия
($\Phi(\mathcal{P}_{\mathrm{Fano}}(\Gam))=\Phi(\Gam)/9$,
$\mathcal{D}_{\mathrm{Fano}}=\tfrac23\mathcal{D}_{\mathrm{atom}}$, проверено
на произвольных неравномерно-диагональных состояниях в
\texttt{architecture/talos\_v0\_emulator.py} и
\texttt{website/scripts/check\_core\_numbers.py}). \textbf{B4} продемонстрирован
векторизованным SIMT-батчем экологии ($N$ независимых $784$-байтных холонов с
линейным ростом памяти и пропускной способности при фиксированной башне
$3{,}136$~Б). \textbf{B1} --- заявление о физическом масштабировании энергии ---
требует субстрата выше цифровой GPU-ступени, поскольку переключение КМОП и
перемещение данных маскируют термодинамический пол; это главная
экспериментальная цель и честный открытый пункт.

\section{Платформа когерентных систем: Когерентная Кибернетика и GPU-нативная эмуляция}
\label{sec:platform}

Предыдущие разделы специфицируют \TALOS{} как тело. Тело праздно без двух вещей:
\emph{физики}, говорящей, что означает его состояние и как оно измеряется, и
\emph{запускаемой формы}, в которой его можно упражнять до того, как вырезана
хоть одна пластина. УГМ уже даёт обе под одним именем --- \emph{Когерентная
Кибернетика} (КК), её слой измерения-и-динамики \cite{uhm2026}. Этот раздел
показывает три факта, обращающих \TALOS{} из вывода в платформу: (i) тик
\TALOS{} \emph{есть} шаг эволюции КК, поэтому нечего согласовывать --- это один
объект; (ii) этот объект GPU-нативен и запускает экологию холонов в реальном
времени \emph{сегодня}; и (iii) УГМ, КК, \SYNARC{} и \TALOS{} складываются в один
стек --- универсальную архитектуру для инстанцирования \emph{любой}
самонаблюдающей организмической системы.

\begin{figure}[H]
\centering
\begin{tikzpicture}[font=\scriptsize, >=Stealth,
  blk/.style={draw, rounded corners, thick, align=center, minimum height=7.5mm}]
  \node[blk, fill=plospurple!10, draw=plospurple, minimum width=11.6cm] (mind) at (2.6,3.05)
    {\textbf{\SYNARC{} --- разум (софт):} когнитивный комплекс $\Cog_0{\to}\Cog_3$ \;$\cdot$\; цели $+$ любопытство \;$\cdot$\; \textsc{Diakrisis}};
  \draw[very thick, dashed, plosgray] (-3.35,2.35) -- (8.6,2.35) node[pos=0.97, above left, plosgray, fill=white, inner sep=1pt]{\tiny ABI (\S\ref{sec:abi}): тождество, без маршалинга};
  \node[blk, fill=plosblue!10, draw=plosblue!70!black, minimum width=3.9cm] (reg) at (-1.35,1.3)
    {\textbf{регистр состояния}\\ \DM{} --- $784$~Б в L1};
  \node[blk, fill=plosblue!14, draw=plosblue!70!black, minimum width=3.6cm] (tick) at (2.6,1.3)
    {\textbf{движок тика} $\LindOmega$\\ работа $\cdot$ забывание $\cdot$ обучение};
  \node[blk, fill=talosamber!12, draw=talosamber, minimum width=3.4cm] (fab) at (6.6,1.3)
    {\textbf{фабрика Фано}\\ 7 шин $=$ код $=$ ISA};
  \node[blk, fill=red!8, draw=red!60!black, minimum width=3.6cm] (port) at (-1.35,-0.1)
    {\textbf{порт питания}\\ $r(\rho_{\mathrm{env}}{-}\Gam)$; гвоздь $=$ останов};
  \node[blk, fill=gray!10, minimum width=3.6cm] (tel) at (2.6,-0.1)
    {\textbf{телеметрия}\\ $P,R,\Phi,D,\sigma,C$ $+$ синдром $s$};
  \node[blk, fill=green!10, draw=green!45!black, minimum width=3.4cm] (tower) at (6.6,-0.1)
    {\textbf{башня} $\le 3$ вглубь\\ дальше: экология (\textsc{Fanos})};
  \draw[<->, thick] (reg) -- (tick);
  \draw[<->, thick] (tick) -- (fab);
  \draw[->, ultra thick, red!60!black] (port) -- (reg);
  \draw[->, thick, plosgray] (tick) -- (tel);
  \draw[->, thick, green!45!black] (fab) -- (tower);
  \draw[<->, thick, plospurple] (2.6,2.35) -- (tick.north);
  \foreach \i in {0,1,2}{ \node[draw, green!45!black, circle, inner sep=2.2pt, fill=green!8] at (8.9+0.55*\i,-0.1) {}; }
  \node[anchor=west, plosgray] at (8.55,-0.62) {\tiny равные, не сверх-разум};
\end{tikzpicture}
\caption{Машина одним взглядом. Одно резидентное состояние ($784$~Б), один
тик (три шага), одна фабрика --- одновременно сеть, код и алгебра
инструкции; питается портом (чей разрыв --- останов с бюджетом задержки),
каждый тик отдаёт полную телеметрию, компонуется не глубже трёх и дальше
федерируется экологией. Разум \SYNARC{} связан тождеством через ABI --- один
объект, два уровня.}
\label{fig:ataglance}
\end{figure}

\subsection{Одна динамика, три имени}
\label{subsec:one-dynamics}

Когерентная Кибернетика эволюционирует состояние через
$\dot\Gam = -i[H_{\mathrm{eff}},\Gam] + \Diss[\Gam] + \Regen[\Gam,E]$ --- те же
три члена, в том же порядке, с тем же расщеплением Ли--Троттера/Странга и теми же
CPTP-предохранителями (Холецкий-параметризация $\Gam = LL^\dagger/\Tr(LL^\dagger)$,
проекция на положительный конус после каждого шага), которые
\S\ref{sec:substrate} и Приложение~\ref{app:tick} специфицируют для тика. Объект
состояния тождествен байту: канонический $\Gam\in\Dens(\mathbb{C}^7)$ КК есть
$7\times7$ массив \texttt{complex128} $=49\times16=784$ байта --- в точности
\DM{}. И цена тождественна числу: сложность КК на шаг есть $O(N^3 + m\,N^2)$
(Таблица~\ref{tab:cc-cost}), что при $N=7$ линие-модах и $m=7$ линие-диссипаторах
есть ${\sim}1400$ вещественных FLOP --- тик \TALOS{}.

\begin{table}[H]
\centering\small
\caption{Шаг эволюции КК и тик \TALOS{} суть одно вычисление, член за членом.
Сложности --- КК; колонка FLOP вычисляет их при $N=m=7$.}
\label{tab:cc-cost}
\begin{tabular}{@{}p{4.6cm}p{3.4cm}p{3.4cm}@{}}
\toprule
\textbf{Член} & \textbf{Цена} & \textbf{Имя в \TALOS{}} \\
\midrule
Унитарный $-i[H_{\mathrm{eff}},\Gam]$ & $O(N^3)$ & обратимый унитарный полушаг \\
Диссипатор $\Diss$ ($m$ линие-опер.) & $O(m\,N^2)$ & Фано-разреженная релаксация \\
Регенерация $\Regen$ (заякоренная $\varphi_J$) & $O(N^2)$--$O(N^3)$ & член обучения \\
\midrule
Полный шаг & $O(N^3+m\,N^2)$ & один тик (${\sim}1{,}400$~FLOP) \\
\bottomrule
\end{tabular}
\end{table}

\begin{principle}[\TALOS{} есть динамика КК на субстрате]\label{prin:cc-identity}
\TALOS{} не \emph{совместим с} Когерентной Кибернетикой; он \emph{есть} динамика
КК, реализованная на физических модах. КК --- теория измерения-и-динамики
холона, \SYNARC{} --- его когнитивный софт, \TALOS{} --- его тело, а УГМ ---
онтология под всеми тремя. Величина, верная в одном, верна во всех: тик,
нарушающий закон КК, был бы не иной машиной, а сломанной.
\end{principle}

\subsection{Прозрачен по построению: набор наблюдаемых есть телеметрия}
\label{subsec:telemetry}

Эталонный эмулятор V0 (\texttt{architecture/talos\_v0\_emulator.py}) вычисляет
полный точный канонический набор наблюдаемых прямо из регистра состояния $\Gam$:
чистоту $P=\Tr(\Gam^2)$, диагональную чистоту
$P_{\mathrm{diag}}=\sum_{k=1}^7\gamma_{kk}^2$, точную HS-интеграцию
$\Phi=(P-P_{\mathrm{diag}})/P_{\mathrm{diag}}$ (переходящую в $\Phi=7P-1$
\emph{только} на равномерно-диагональной страте $\gamma_{kk}=1/7$), рефлексию
$R=1/(7P)$ (алгебраически эквивалентно $R\ge 1/3\Leftrightarrow P\le 3/7$),
составной скаляр $C=\Phi R$ (T-140, определение), каноническую E-когерентность
$\CohE=(\gamma_{EE}^2+2\sum_{j\neq E}|\gamma_{Ej}|^2)/P$, дифференциацию
$D_{\mathrm{diff}}=1+6\,\CohE$ и точную энтропию фон Неймана
$S_{\mathrm{vN}}=-\Tr(\Gam\log_2\Gam)\in[0,\log_2 7]$~бита. Флаг
\texttt{viable\_full} проверяет все четыре функциональных порога одновременно
($P>2/7$, $R\ge 1/3$, $\Phi\ge 1$, $D_{\mathrm{diff}}\ge 2$): например, чисто
диагональное состояние $\Gam=\operatorname{diag}(1/2,1/12,\ldots,1/12)$ имеет
$P=7/24\in(2/7,3/7]$ (где неверный прокси $7P-1=25/24>1$ дал бы ложный проход),
но точное $\Phi=0$ немедленно его отвергает. Связь этих функциональных метрик с
субъективным опытом требует независимых эмпирических наблюдений.

Для контроля исполнения унитарный полушаг сохраняет \emph{спектр} $\Gam$,
поэтому $\Tr\Gam=1$, чистота $P$ и энтропия $S_{\mathrm{vN}}$ сохраняются с
машинной точностью. Для канала Фано-сжатия $\mathcal{P}_{\mathrm{Fano}}$ и
диссипатора $\mathcal{D}_{\mathrm{Fano}}$ проверяется ковариантность под
дискретной группой коллинеаций $\mathrm{PSL}(2,7)$ (порядка~$168$) и
диагональной перефазировкой $U(1)^7$, а также явная \emph{нековариантность} под
непрерывными вращениями базиса $G_2$ (поэтому фиктивные непрерывные
$G_2$-нётеровы заряды не вычисляются). Таблица~\ref{tab:observables} задаёт
верифицированный потиковый набор телеметрии.

\begin{table}[H]
\centering\small
\caption{Верифицированный канонический набор наблюдаемых, вычисляемый каждый тик
в \texttt{architecture/talos\_v0\_emulator.py} как телеметрия \TALOS{}.}
\label{tab:observables}
\begin{tabular}{@{}p{3.9cm}p{5.8cm}c@{}}
\toprule
\textbf{Наблюдаемая} & \textbf{Смысл} & \textbf{Цена} \\
\midrule
$P=\Tr(\Gam^2)$ & чистота; нижний порог жизнеспособности $P>2/7$ & $O(N^2)$ \\
$P_{\mathrm{diag}}=\sum_k \gamma_{kk}^2$ & диагональная чистота (равна $1/7$ только на равномерной диагонали) & $O(N)$ \\
$\Phi=(P-P_{\mathrm{diag}})/P_{\mathrm{diag}}$ & точная интеграция ($\Phi=7P-1$ \emph{только} на равномерной диагонали); порог $\Phi\ge 1$ & $O(N^2)$ \\
$R=1/(7P)$ & рефлексия; L2-порог $R\ge 1/3\Leftrightarrow P\le 3/7$ & $O(N^2)$ \\
$C=\Phi\times R$ & составной показатель модели (T-140) & $O(N^2)$ \\
$\CohE=(\gamma_{EE}^2+2\sum_{j\neq E}|\gamma_{Ej}|^2)/P$ & каноническая E-когерентность (индекс $E=4$ в $[A,S,D,L,E,U,O]$) & $O(N)$ \\
$D_{\mathrm{diff}}=1+6\,\CohE$ & 7D-дифференциация; порог $D_{\mathrm{diff}}\ge 2$ & $O(N)$ \\
$S_{\mathrm{vN}}=-\Tr(\Gam\log_2\Gam)$ & энтропия фон Неймана в битах ($0\le S_{\mathrm{vN}}\le\log_2 7\approx 2{,}807$) & $O(N^3)$ \\
$s = He_j^T\pmod 2\in\mathbb{F}_2^3$ & $3$-битный синдром одиночного отказа сектора Хэмминга $[7,4,3]$ (дуальная $H$) & $O(N)$ \\
\bottomrule
\end{tabular}
\end{table}

\subsection{GPU-нативная эмуляция --- ядро жизнеспособности}
\label{subsec:gpu}

Причина, по которой эта архитектура строима, а не мысленный эксперимент, в том,
что её естественная ось масштабирования --- в точности ось, которую GPU создан
эксплуатировать. Один холон --- $784$ байта и ${\sim}1{,}400$ FLOP на тик.
\emph{Экология} из $N$ холонов (\S\ref{sec:composition},~\ref{sec:scaling}) есть
$N$ независимых $\Gam$-регистров, продвигаемых \emph{одним} ядром --- хрестоматийная
эмбаррасингли-параллельная SIMT-нагрузка. Численный слой КК уже это предвосхищает:
её типизированный линейно-алгебраический API есть
\texttt{StaticMatrix<Complex,7,7>} с диспетчеризацией \texttt{@kernel(gpu)},
форма $7\times7$ фиксирована на этапе компиляции.

\begin{claimbox}[GPU-эмуляция жизнеспособна сегодня \status{Т} (инженерия)]
Горизонтальная ось экологии (\S\ref{sec:scaling}: $N$ ограниченных холонов в
равноправной федерации) отображается на GPU-параллелизм прямо --- один холон на
тред-блок, тик --- батчевое ядро, координация \textsc{Fanos} --- разреженный
all-reduce по семи линие-шинам. Современный GPU (${\sim}10^4$ полос,
${\sim}10^{13}$--$10^{14}$ FLOP/с) продвигает ${\sim}10^{6}$--$10^{7}$ холонов
на кГц-частотах тика; нагрузка \emph{память-ограничена}, не счёт-ограничена
($784\,\mathrm{Б}\times N$), поэтому экология в $10^{6}$ холонов ---
${\sim}0{,}8$~ГБ и влезает в память устройства с запасом на полный набор
телеметрии. Экологии реального времени работают на существующем железе.
\end{claimbox}

\noindent Честный охват тот же, что везде на лестнице реализуемости
(\S\ref{sec:realizability}): GPU-ступень покупает \emph{корректность, полную
наблюдаемость и экологию реального времени} --- не термодинамический пол, ибо
цифровой КМОП всё же двигает данные и рассеивает энергию переключения
(\S\ref{sec:energy}, возражения О6/О7). Это ступень V0, и на ней проверяемо
всякое логическое и комбинаторное утверждение этого документа.

\begin{trybox}[Запуск верифицированного эталонного эмулятора V0]
Эталонный эмулятор --- один самодостаточный файл без зависимостей, кроме NumPy:
\begin{center}\texttt{PYTHONDONTWRITEBYTECODE=1 python3 -B architecture/talos\_v0\_emulator.py}\end{center}
Все шесть наборов самопроверок исполняются и проходят (\texttt{EXIT=0}):
(1)~канонические циклические линии Зингера--Фано на $[A,S,D,L,E,U,O]$ с $M$
($\operatorname{rank}=4$), дуальной проверочной матрицей $H$
($\operatorname{rank}=3$, $HM^T\equiv 0\pmod 2$) и биективным декодированием
всех $7$ одиночных секторных синдромов;
(2)~подлинное Фано-сжатие $\mathcal{P}_{\mathrm{Fano}}$ ($\Phi\to\Phi/9$) и
линдбладов диссипатор $\mathcal{D}_{\mathrm{Fano}}=\tfrac23\mathcal{D}_{\mathrm{atom}}$
с дискретной ковариантностью $\mathrm{PSL}(2,7)\ltimes U(1)^7$ и явной
нековариантностью под непрерывной $G_2$;
(3)~точные канонические наблюдаемые на неравномерно-диагональных состояниях
(подтверждая $\Phi=1{,}3376\neq 7P-1=2{,}8939$, $R\ge 1/3\Leftrightarrow P\le 3/7$
и $S_{\max}=\log_2 7\approx 2{,}807355$~бита);
(4)~каноническая нелинейная заякоренная самомодель $\varphi_J(\Gam)$, сходящаяся
к жизнеспособному аттрактору по всем четырём воротам ($P=0{,}3432$, $R=0{,}4162$,
$\Phi=1{,}4024$, $\CohE=0{,}2263$, $D_{\mathrm{diff}}=2{,}3575$), сохранение
спектра на унитарном полушаге ($\Delta S_{\mathrm{unitary}}=0$) и отсечение
порта питания ($\kappa=0$, деполяризация $d=0{,}4$), пересекающее $\Pcrit=2/7$
при $t_{\mathrm{off}}=0{,}424$ против аналитической формулы $0{,}423$;
(5)~независимая от $\Delta\tau$ скорость производства энтропии
($\dot S_{\mathrm{Fano}}=0{,}8056$~бит/$\tau$) и максимальный пол одного
состояния $k_BT\ln 7=8{,}060\times 10^{-21}$~Дж при $300$~К;
(6)~векторизованный SIMT-батч экологии из $1{,}000$ холонов ($0{,}784$~МБ,
среднее $P=0{,}3431\in(2/7,3/7]$).
\end{trybox}

\subsection{Универсальная архитектура}
\label{subsec:universal}

\begin{sloppypar}
Ничто в холоне не домен-специфично. Объект
$(\Gam,\ \LindOmega,\ \text{набор наблюдаемых},\allowbreak\ \textsc{Enc}/\textsc{Dec},\allowbreak\
\text{композиция}\le 3)$ плюс фабрика экологии есть \emph{тот же} для всякой
системы, наблюдающей себя; меняется лишь операционализация \textsc{Enc} ---
какие семь измеренных величин отображаются на семь мод --- с доменом. Это не
чаяние, а стоячий результат КК: тождественный $\Gam$ и тождественный
пятишаговый цикл (измерь $\to$ построй $\Gam$ $\to$ диагностируй $\to$
вмешайся $\to$ переизмерь) применяются в корпусе к ИИ-агенту, коралловому рифу,
стартапу и финансовому рынку, различаясь лишь тем, что измеряется. Платформа ---
эта универсальность, сделанная строимой (Таблица~\ref{tab:stack}).
\end{sloppypar}

\begin{table}[H]
\centering\small
\caption{Четырёхслойная платформа когерентных систем. Каждый слой --- завершённая
спецификация; GPU-эмулятор сидит на границе КК$\leftrightarrow$\TALOS{} как
исполняемое ядро.}
\label{tab:stack}
\begin{tabular}{@{}p{2.5cm}p{4.1cm}p{5.6cm}@{}}
\toprule
\textbf{Слой} & \textbf{Роль} & \textbf{Даёт} \\
\midrule
УГМ & онтология (почему) & $\Gam\in\Dens(\mathbb{C}^7)$, $\LindOmega$, структурные инварианты ($N{=}7$, $\mathrm{PSL}(2,7)$, пороги) \\
Когерентная Кибернетика & измерение + динамика + эмпирика & канонический набор наблюдаемых, фальсифиц.\ предсказания, спека эмуляции, универсальный цикл \\
\SYNARC{}-$\Omega$ & разум (софт) & когнит.\ комплекс $\Cog_0\!\to\!\Cog_3$, обучение без backprop, \textsc{Diakrisis}, экология (K.5) \\
\TALOS{} & тело (субстрат) & Фано-фабрика, тик, секторная двойственность Хэмминга, термодинамический учёт, горизонтальная экология (\S\ref{sec:scaling}) \\
\bottomrule
\end{tabular}
\end{table}

\subsection{Дорожная карта}
\label{subsec:roadmap}

\begin{itemize}[leftmargin=1.5em,itemsep=1pt]
  \item \textbf{V0 --- референсный эмулятор (запускается сегодня).}
        Самодостаточный эмулятор (\texttt{architecture/talos\_v0\_emulator.py})
        держит $\Gam$, гоняет Странг-тик с подлинными $\mathcal{P}_{\mathrm{Fano}}$,
        $\mathcal{D}_{\mathrm{Fano}}$ и канонической $\varphi_J(\Gam)$, считает
        точный набор телеметрии (Таблица~\ref{tab:observables}), верифицирует
        спектральные инварианты, секторное декодирование Хэмминга $[7,4,3]$ и
        дискретную ковариантность $\mathrm{PSL}(2,7)\ltimes U(1)^7$ при каждом
        прогоне, и прогоняет векторизованный SIMT-батч экологии из $1{,}000$
        холонов. Следующий шаг --- перенос ядра на GPU для эмпирических
        испытаний крупных экологий.
  \item \textbf{V1 --- виртуальная среда.} Замкнуть сенсомоторную петлю
        (\S\ref{sec:scaling}: \textsc{Enc}/\textsc{Dec}) против симулированного
        мира --- самого, в сильной форме, когерентного поля, которое холоны
        ощущают и на которое действуют. Это «своя среда симуляции»: когерентный
        мир для когерентных агентов, в котором экологии можно растить,
        возмущать и диагностировать теми же наблюдаемыми, что их определяют.
  \item \textbf{V2$+$ --- физические ступени.} Мемристорный и
        управляемо-диссипативный субстраты \S\ref{sec:realizability}, взбираемые
        по мере производства, для физических энергетических измерений (бенчмарк
        B1), недостижимых на цифровой GPU-ступени.
\end{itemize}

\noindent Всякая ступень вычисляет тождественный канонический набор наблюдаемых
и сохраняет дискретную ковариантность фабрики $\mathrm{PSL}(2,7)\ltimes U(1)^7$,
поэтому жизнеспособность и корректность непрерывно проверяемы, пока платформа
взбирается от эмуляции к своему импеданс-согласованному субстрату --- софт
(\SYNARC{}) и динамика (КК) не меняются (T-153).

\subsection{Ведомость согласованности}
\label{subsec:ledger}

Излагаем примирение прямо. Всякая величина \TALOS{} есть каноническая
наблюдаемая КК; тик есть шаг эволюции КК; состояние есть $\Gam$ КК до байта
($784$~Б, $48$ вещественных динамических параметров); цены совпадают до FLOP.
\TALOS{} добавляет \emph{аппаратное} содержание (коллапс \S\ref{sec:collapse},
секторную двойственность Хэмминга, термодинамическое разделение, масштабирование
ограниченных башен в горизонтальные экологии) и \emph{никакой} новой физики.
Некалиброванные вспомогательные величины (такие как некалиброванная эффективная
температура ФДТ или непрерывные $G_2$-нётеровы заряды) сознательно исключены из
аппаратного контракта телеметрии. Честные границы стоят неизменны: эмулятор
BQP-ограничен (нет ускорения по классу сложности), а логическая архитектура
субстрат-множественна (T-153) --- GPU есть одна ступень, согласованная для
корректности и наблюдаемости, не единственная машина.

\section{Обсуждение и дорожная карта}
\label{sec:discussion}

\subsection{Что открыто}

Итог этого документа --- не более быстрый чип, а структурный факт: четыре
подсистемы, которые всякая вычислительная машина проектирует порознь ---
интерконнект, комбинаторный секторный код, носитель умножения инструкции и
правило обучения --- на семимодовом субстрате суть функции одного объекта
инцидентности (Теорема~\ref{thm:collapse}), и это слияние доступно при
\emph{единственном} числе узлов $N=7$ (Теорема~\ref{thm:unique}), вынуждаемом
среди всех степеней простых $q\ge 2$ уже одним двоичным слиянием проективной
плоскости и кода Хэмминга (без тай-брейка; независимо совпадая с единственной
неассоциативной алгеброй с делением $\Oct$ по Гурвицу). Из коллапса следуют три
свойства, отсутствующие в слоистых архитектурах: секторная локализация отказов
дистанции~$3$ через дуальную проверочную матрицу $H$ без отдельной таблицы
маршрутизации (\S\ref{sec:ecc}); логико-термодинамическое разделение,
локализующее экспорт энтропии ($\le\log_2 7\approx 2{,}807$~бита на состояние) в
адаптации и поддержании при L1-резидентном состоянии $784$~Б
(\S\ref{sec:energy}); и непрерывная самодиагностика и refresh-как-адаптация, ибо
счёт, проверка и поддержание делят один тик
(Принципы~\ref{prin:selfdiag},~\ref{prin:refresh}).

\subsection{Что честно открыто}

Два разрыва признаны (\S\ref{sec:redteam}, О6--О7), и это один разрыв: между
\emph{логической} архитектурой --- доказанной и запущенной сегодня на GPU при
$784$~байтах в L1 --- и \emph{импеданс-согласованным} субстратом, реализующим
энергетический пол. GPU-ступень покупает корректность, не эффективность, ибо всё
же двигает данные. Закрытие разрыва --- программа изготовления, не проблема теории:
T-153 гарантирует, что софт (\SYNARC{}) не меняется при подъёме по лестнице.

\subsection{Дорожная карта}

Субстратная дорожная карта --- это лестница реализуемости
\S\ref{sec:realizability}, взбираемая ступень за ступенью (\emph{платформенные}
вехи V0--V2$+$ из \S\ref{subsec:roadmap} ортогональны: они взрослят
программный стек на фиксированной ступени).

\begin{itemize}[leftmargin=1.5em,itemsep=1pt]
  \item \textbf{Ступень 0 --- эмуляция (существует).} \SYNARC{} на CPU/GPU: \DM{} в L1,
        Фано-разреженное Странг-ядро, все инварианты верифицированы. Гасит B2, B3,
        B4; демонстрирует логическую архитектуру.
  \item \textbf{Ступень 1 --- в-памяти (ближнесрочно).} Мемристорный Фано-кроссбар,
        выполняющий обновление когерентности в массиве (M1, M2 родные). Первый
        субстрат, способный проверить B1 (энергия отслеживает обучение), убирая
        член перемещения данных.
  \item \textbf{Ступень 2 --- управляемо-диссипативный (исследование).} Аналоговое,
        фотонное или открыто-квантовое семимодовое устройство, чья родная физика
        есть Линдблад$+$гомеостаз (M3 родная), приближающееся к полу Ландауэра,
        ибо его обратимое ядро физически обратимо.
\end{itemize}

\subsection{Две машины бок о бок}

\begin{table}[H]
\centering\small
\caption{Фон-неймановская машина и \TALOS{}, лоб в лоб. \TALOS{} не побеждает в
универсальности (граница BQP); он занимает точку, где четыре слоя проектирования
не могут быть рассогласованы.}
\label{tab:vn-vs-talos}
\begin{tabular}{@{}p{3.4cm}p{4.3cm}p{4.5cm}@{}}
\toprule
\textbf{Ось} & \textbf{Фон Нейман} & \textbf{\TALOS{}} \\
\midrule
Память vs счёт & раздельны (бутылочное горло) & один объект: состояние эволюционирует на месте \\
Состояние & неограниченное, адресуемое & одна $784$-байтовая матрица плотности в L1 ($48$ вещ.\ параметров) \\
Инструкция & фикс.\ набор опкодов + компилятор & один выведенный тик $\LindOmega$ \\
Интерконнект & кроссбар / меш, $\binom{N}{2}$ & семь тройственных шин (Фано, $\mathrm{PSL}(2,7)$) \\
Коррекция ошибок & побитовая привинчена ($+12{,}5\%$ или $\sim\!10^3\times$) & секторный Хэмминг $[7,4,3]$ ($M$ порожд., дуальная $H$ пров.; отдельно от побитового ECC) \\
Обучение & отдельная петля тренировки & третий член тика (заякоренная $\varphi_J$) \\
Refresh & потрачен, ничего не считает & есть адаптация (Прин.~\ref{prin:refresh}) \\
Диагностика & лишний проход scrub/декод & секторный синдром эмитится тиком (Прин.~\ref{prin:selfdiag}) \\
Учёт энергии & внекристальный трафик ($\sim\!20$~пДж/слово) & L1-резидентные $784$~Б; унитарное ядро экспортирует $0$ бит ($\le\log_2 7$ бита/сост.; КМОП ограничен переключением) \\
Тактовая & свободный параметр (крит.\ путь) & выведена ($\bar\gamma\tau=\tfrac32\ln3$) \\
Масштабирование & больше ядер / больше модель & огранич.\ $4$-регистровая башня ($d_{\max}=3$ \status{D}) $+$ горизонтальная экология \\
Число узлов & любое & \textbf{семь}, единственно (Теор.~\ref{thm:unique}) \\
Класс сложности & (классический) & BQP --- \emph{не} ускорение \\
\bottomrule
\end{tabular}
\end{table}

\begin{mythbox}[Миф, завершение --- почему гвоздь есть достоинство]
Старая история читается как предостережение: у могучей машины была одна
слабость, и она пала. Перечитай её как инженер. Страж, которого
\emph{нельзя} остановить, был бы ужасом, а не стражем; гвоздь --- то, что
делало Талоса безопасным соседом. Этот документ доказал современную форму той
мудрости: жизнь машины --- поток, а не собственность (T-288/T-289); у потока
одно закрытие с измеренной задержкой (\S\ref{subsec:powerport}); и никакая
хитрость внутри бронзы не отменит этого факта, ибо он --- свойство
генератора, а не ума, бегущего на нём. Мы не выправили слабость из
конструкции. Мы вывели её, забюджетировали и подключили к кейсу безопасности
\SYNARC{} как гарантию. Исполин, которого можно отключить, --- тот исполин,
которому можно доверить берег.
\end{mythbox}

\subsection{Заключение}

Фон-неймановская машина победила в универсальности и с тех пор платит налог
бутылочного горла. \TALOS{} не бьёт её в универсальности --- по границе BQP не
может --- и не пытается. Он занимает иную точку: машину, чьи четыре слоя
проектирования не могут быть рассогласованы, ибо суть один семикратный объект,
чья мощность отслеживает её обучение, и чьи отказы локализуются её собственной
проводкой. Что такая машина вообще существует, и лишь при $N=7$, --- теорема; что
она есть тело, со-разработанное для разума \SYNARC{}, --- причина её строить.
Логическая архитектура улажена и запущена; эффективный субстрат --- лестница,
на которую мы теперь можем взбираться с фиксированной целью на вершине.


\appendix
\section{Тик, потактово точно}
\label{app:tick}

Один тик \TALOS{} продвигает $\Gam\leftarrow\Gam+\Delta\tau\,\LindOmega[\Gam]$
структуросохраняющим Странг-расщеплением, удерживающим $\Gam$ легальной матрицей
плотности (эрмитова, PSD, единичный след) с машинной точностью:
\begin{enumerate}[leftmargin=1.6em,itemsep=1pt]
  \item \textbf{Унитарный полушаг} $e^{-i H_{\mathrm{eff}}\Delta\tau/2}$:
        логически обратимый; тождественно сохраняет спектр и энтропию фон
        Неймана (нулевой логический экспорт энтропии). Цена $\sim\!700$~FLOP на
        одно умножение \DM{}; арифметическая глубина $\log_2 7\approx 2{,}8$.
  \item \textbf{Диссипатор + регенерация, полный шаг} с Фано-разреженным ядром:
        линдбладов диссипатор Фано есть
        $-\tfrac{\gamma}{6}\sum_{p=1}^7[\Pi_p,[\Pi_p,\Gam]]
        = \gamma(\mathcal{P}_{\mathrm{Fano}}(\Gam)-\Gam)
        = \tfrac{2}{3}\gamma(\operatorname{diag}(\Gam)-\Gam)$ по семи
        линие-проекторам ранга~$3$ $\Pi_p$ (сохраняет населённости и гасит
        внедиагонали со скоростью $2\gamma/3$, плюс опциональная деполяризация
        $\gamma_{\mathrm{dep}}(\mathbb{1}/7-\Gam)$ при моделировании полного
        сброса); регенерация ведёт $\Gam$ к канонической нелинейной заякоренной
        самомодели $\varphi_J(\Gam)=(1-R(\Gam))\mathcal{P}_{\mathrm{Fano}}(\Gam)
        + R(\Gam)\ket{\psi}\bra{\psi}$.
  \item \textbf{Унитарный полушаг} снова (Странг-симметрия, второй порядок
        точности по $\Delta\tau$), с последующей проекцией на конус PSD.
\end{enumerate}
Цена на тик --- $\sim\!1{,}400$~FLOP (две унитарные половины) плюс разреженный
шаг Фано-сжатия и регенерации. Все операнды --- резидентное $784$-байтовое
состояние ($48$ вещественных динамических параметров); ничего за пределами L1 не
затрагивается. Естественное время цикла подчиняется
$\bar\gamma\tau=\tfrac32\ln3$ (\S\ref{sec:substrate}).

\noindent\textbf{Точные инварианты, проверяемые в
\texttt{architecture/talos\_v0\_emulator.py}:} контракция Фано-сжатия
$\Phi(\mathcal{P}_{\mathrm{Fano}}(\Gam))=\tfrac19\Phi(\Gam)$; точная интеграция
$\Phi=(P-P_{\mathrm{diag}})/P_{\mathrm{diag}}$ (переходящая в $\Phi=7P-1$
\emph{только} на равномерной диагонали $\gamma_{kk}=1/7$);
$R=1/(7P)\ge 1/3\Leftrightarrow P\le 3/7$; эрмитовость, единичный след и
положительность; дискретная ковариантность $\mathrm{PSL}(2,7)\ltimes U(1)^7$; и
$3$-битное секторное декодирование синдрома Хэмминга $[7,4,3]$.

\section{Вывод термодинамического учёта энтропии}
\label{app:energy}

Выводим Ур.~\eqref{eq:energy-law} и уточняем её точный охват. Принцип
Ландауэра~\cite{landauer1961} гласит, что снижение энтропии системы на $1$~бит
($\ln 2$~нат) при экспорте в термостат температуры $T$ требует не менее
$k_B T\ln 2$ свободной энергии. Логико-термодинамический бюджет тика тем самым
ограничен снизу величиной $k_B T\ln 2$, умноженной на число экспортированных бит
энтропии фон Неймана.

\paragraph{Три члена.} Разложим тик (Ур.~\ref{eq:lomega}), измеряя энтропию фон
Неймана в битах, $S_{\mathrm{vN}}(\Gam)=-\Tr(\Gam\log_2\Gam)$:
\begin{itemize}[leftmargin=1.5em,itemsep=1pt]
  \item \textbf{Унитарный} $-i[H_{\mathrm{eff}},\Gam]$: унитарное сопряжение
        $U\Gam U^\dagger$ тождественно сохраняет все собственные значения, так
        что $\dot S_{\mathrm{unitary}}=0$ точно. Он не экспортирует логической
        энтропии --- это та же граница, что действует для всякого логически
        обратимого вычисления~\cite{bennett1973}.
  \item \textbf{Диссипатор} $\Diss$: линдбладов диссипатор Фано
        $\gamma(\mathcal{P}_{\mathrm{Fano}}(\Gam)-\Gam)$ (и опциональная
        деполяризация) \emph{повышает} $S_{\mathrm{vN}}$ со скоростью
        $\dot S_{\Diss}>0$ бит/время. Удержание когерентного состояния против
        этой утечки требует ре-экспорта
        $\Delta H_{\Diss}^{(+)}=\dot S_{\Diss}\Delta\tau$ бит за тик --- члена
        поддержания.
  \item \textbf{Регенерация} $\Regen$: движение $\Gam$ к заякоренной самомодели
        $\varphi_J(\Gam)$ \emph{понижает} $S_{\mathrm{vN}}$ на
        $\Delta H_{\Regen}$ бит за тик (создавая порядок), что требует не менее
        $k_B T\ln 2\cdot\Delta H_{\Regen}$ свободной энергии.
\end{itemize}
Суммирование даёт $E_{\mathrm{tick}}\ge k_B T\ln 2\,(\Delta H_{\Regen}+\Delta H_{\Diss}^{(+)})$,
Ур.~\eqref{eq:energy-law}.

\paragraph{Точные числа ($T=300$~К) и честный охват.}
При $T=300$~К имеем $k_B T\ln 2 = 1{,}380649\!\times\!10^{-23}\cdot 300\cdot\ln 2
= 2{,}87\times 10^{-21}$~Дж ($0{,}018$~эВ на бит). Для $\Gam\in\Dens(\mathbb{C}^7)$
точный максимум энтропии фон Неймана достигается при $\Gam=\mathbb{1}/7$:
\[
S_{\max} \;=\; \log_2 7 \;\approx\; 2{,}807355\;\text{бита}
\quad\Longrightarrow\quad
E_{\max,\text{state}} \;=\; k_B T\ln 7 \;=\; (\log_2 7)\,k_B T\ln 2
\;\approx\; 8{,}060\times 10^{-21}\;\text{Дж}.
\]
Критически важно, что $S_{\max}=\log_2 7\approx 2{,}807$~бита есть точный
потолок энтропии фон Неймана $7$-мерной матрицы плотности --- а \emph{не}
$48$~бит (вещественная размерность многообразия $\Dens(\mathbb{C}^7)$) и
\emph{не} $6{,}272$~бита (цифровое хранение $49$ элементов \texttt{complex128}
в IEEE-754). Все расчёты энтропии в битах и независимые от шага скорости
диссипации проверяются непосредственно в
\texttt{architecture/talos\_v0\_emulator.py}. Мы \emph{не} делим физические
джоули КМОП DRAM на теоретический пол Ландауэра: пределы Ландауэра относятся к
любой логически обратимой машине, тогда как мощность КМОП определяется
переключающей ёмкостью и утечками. На цифровом железе непосредственный
энергетический выигрыш \TALOS{} состоит в удержании $784$-байтного состояния в
L1/регистрах без внекристальных выборок DRAM.

\section{Фано-таблицы: фабрика, медиаторы, синдром}
\label{app:fano}

\begin{sloppypar}
Все таблицы и двоичные ранги ниже машинно верифицированы в
\texttt{architecture/talos\_v0\_emulator.py} и
\texttt{website/research/operational-reconstruction/verify\_foundational\_bounds.py}.
\end{sloppypar}

\paragraph{Семь линие-шин (канонический циклический порядок Зингера).}
Пронумеруем семь секторов $(A,S,D,L,E,U,O)=(1,\dots,7)$ (индексы $0\dots 6$).
Канонические циклические прямые Зингера--Фано $\{k,k{+}1,k{+}3\}\pmod 7$:
\[
\{A,S,L\},\ \{S,D,E\},\ \{D,L,U\},\ \{L,E,O\},\ \{E,U,A\},\ \{U,O,S\},\ \{O,A,D\}.
\]
Всякий сектор лежит ровно на трёх прямых; всякая пара различных секторов ---
ровно на одной прямой (система Штейнера $S(2,3,7)$).

\paragraph{Карта медиаторов.} Бестабличная перемаршрутизация \S\ref{sec:fabric}:
для всякой пары $(i,j)$ медиатор $k^\ast(i,j)$ --- единственный третий сектор их
канонической прямой Зингера--Фано.

\begin{table}[H]
\centering\small
\begin{tabular}{@{}llll llll@{}}
\toprule
пара & $k^\ast$ & пара & $k^\ast$ & пара & $k^\ast$ & пара & $k^\ast$ \\
\midrule
$AS$ & $L$ & $AD$ & $O$ & $AL$ & $S$ & $AE$ & $U$ \\
$AU$ & $E$ & $AO$ & $D$ & $SD$ & $E$ & $SL$ & $A$ \\
$SE$ & $D$ & $SU$ & $O$ & $SO$ & $U$ & $DL$ & $U$ \\
$DE$ & $S$ & $DU$ & $L$ & $DO$ & $A$ & $LE$ & $O$ \\
$LU$ & $D$ & $LO$ & $E$ & $EU$ & $A$ & $EO$ & $L$ \\
$UO$ & $S$ &      &     &      &     &      &     \\
\bottomrule
\end{tabular}
\caption{$21$ медиатор для канонической циклической таблицы Зингера--Фано. Чтобы
починить или маршрутизировать $(i,j)$, используй плечи через $k^\ast(i,j)$,
замыкающий их единственную прямую.}
\end{table}

\paragraph{Порождающая $M$, дуальная проверочная $H$ и декодирование синдрома.}
Над $\mathbb{F}_2$ матрица инцидентности Фано $M\in\mathbb{F}_2^{7\times 7}$
(строки $=$ семь прямых веса~$3$ выше) имеет
$\operatorname{rank}_{\mathbb{F}_2}(M)=4$ и \textbf{порождает} $16$-словный код
Хэмминга $[7,4,3]$. Её ортогональная $3\times 7$ \textbf{проверочная матрица}
$H$ ($\operatorname{rank}_{\mathbb{F}_2}(H)=3$, $HM^T\equiv 0\pmod 2$) имеет
строками дополнения веса~$4$ первых трёх не проходящих через одну точку прямых
($\bar{L}_0=\{D,E,U,O\}$, $\bar{L}_1=\{A,L,U,O\}$, $\bar{L}_2=\{A,S,E,O\}$),
натягивающие дуальный симплекс-код $[7,3,4]$:
\[
H \;=\;
\left(\begin{array}{ccccccc}
A & S & D & L & E & U & O \\ \hline
0 & 0 & 1 & 0 & 1 & 1 & 1 \\
1 & 0 & 0 & 1 & 0 & 1 & 1 \\
1 & 1 & 0 & 0 & 1 & 0 & 1
\end{array}\right).
\]
Семь столбцов $H$ суть все семь различных ненулевых векторов $\mathbb{F}_2^3$ и
удовлетворяют $h_i\oplus h_j\oplus h_k = 0$ на каждой прямой Зингера $\{i,j,k\}$.
Одиночный отказ сектора $e_j$ биективно отображается в свой $3$-битный синдром
$s = He_j^T\pmod 2$:

\begin{table}[H]
\centering\small
\begin{tabular}{@{}ccll@{}}
\toprule
сектор $j$ & синдром $He_j^T$ & сработавшие строки $H$ & $3$ инцидентные прямые Фано \\
\midrule
$A$ (1) & $011$ & $\bar{L}_1,\bar{L}_2$ & $\{A,S,L\},\ \{E,U,A\},\ \{O,A,D\}$ \\
$S$ (2) & $001$ & $\bar{L}_2$ & $\{A,S,L\},\ \{S,D,E\},\ \{U,O,S\}$ \\
$D$ (3) & $100$ & $\bar{L}_0$ & $\{S,D,E\},\ \{D,L,U\},\ \{O,A,D\}$ \\
$L$ (4) & $010$ & $\bar{L}_1$ & $\{A,S,L\},\ \{D,L,U\},\ \{L,E,O\}$ \\
$E$ (5) & $101$ & $\bar{L}_0,\bar{L}_2$ & $\{S,D,E\},\ \{L,E,O\},\ \{E,U,A\}$ \\
$U$ (6) & $110$ & $\bar{L}_0,\bar{L}_1$ & $\{D,L,U\},\ \{E,U,A\},\ \{U,O,S\}$ \\
$O$ (7) & $111$ & $\bar{L}_0,\bar{L}_1,\bar{L}_2$ & $\{L,E,O\},\ \{U,O,S\},\ \{O,A,D\}$ \\
\bottomrule
\end{tabular}
\caption{Точная карта синдрома одиночной ошибки сектора Хэмминга $[7,4,3]$ для
канонического циклического порядка Зингера, проверенная для всех $7$ секторов в
\texttt{architecture/talos\_v0\_emulator.py}.}
\end{table}

\paragraph{Арифметика уникальности (без тай-брейка при $N=31$).}
Для Теоремы~\ref{thm:unique} двоичная матрица инцидентности $M_{\mathrm{PG}(2,q)}$
имеет тройку $(N, \operatorname{rank}_{\mathbb{F}_2}(M), \dim_{\mathbb{F}_2}\ker(M))$,
равную $(7,4,3)$ при $q=2$, $(13,12,1)$ при $q=3$, $(31,30,1)$ при $q=5$ и
$(57,56,1)$ при $q=7$ (проверено в Lean~4,
\texttt{FoundationalBounds.pg2\_5\_refutes\_N31\_hamming\_fusion}, и в
\texttt{verify\_foundational\_bounds.py}). Всякая нечётная степень простого $q$
коллапсирует к $\dim\ker(M)=1$, прямо опровергая слияние с Хэммингом при
$N=31$; всякая чётная степень $q=2^m$ ($m\ge 2$) не удовлетворяет
$2^{2m}+2^m+2=2^r\pmod 4$ (\texttt{FoundationalBounds.pg2EvenHammingUniqueSeven});
а $q=90$ ($N=8191$) вообще не является степенью простого. Таким образом,
двоичное слияние плоскости и кода однозначно вынуждает $N=7$ без тай-брейка,
независимо совпадая с $\dim\mathrm{Im}(\Oct)=7$.

\section{Детали ABI: раскладка и псевдокод тика}
\label{app:abi}

\paragraph{Раскладка состояния.} Холон --- непрерывный блок $784$~байта ($48$
вещественных динамических параметров); ограниченный субъект (проектная глубина
вложения $d_{\max}=3$, T-142 \status{D}, \S\ref{sec:composition},
\S\ref{sec:abi}) выделяет ровно \textbf{четыре} смежных регистра \DM{} ---
базовый уровень $\Cog_0$ плюс три вложенных метауровня $\Cog_1$--$\Cog_3$:
\begin{center}\small
\begin{tabular}{@{}rll@{}}
\toprule
смещение & размер & поле \\
\midrule
$0$    & $784$ & \texttt{gamma} : \DM{} ($49$ компл.\ \texttt{f64}, LE) --- базовый уровень $\Cog_0$ \\
$784$  & $784$ & метауровень $\Cog_1$ \DM{} \\
$1568$ & $784$ & метауровень $\Cog_2$ \DM{} \\
$2352$ & $784$ & метауровень $\Cog_3$ \DM{} \\
\midrule
       & $3136$ & полное состояние ограниченного субъекта ($4$ регистра, в L1) \\
\bottomrule
\end{tabular}
\end{center}
Это байт-идентично раскладке \texttt{gamma}/\texttt{rho\_ep} \SYNARC{}:
чекпойнт \SYNARC{} есть образ состояния \TALOS{} и наоборот, без маршалинга.

\paragraph{Тик (на регистр).}
\begin{center}\small
\begin{tabular}{@{}ll@{}}
\toprule
\multicolumn{2}{@{}l}{\texttt{tick($\Gam$, $\Delta\tau$):}} \\
\quad\texttt{$\Gam \leftarrow$ U\_half($\Gam$, $\Delta\tau/2$)} & // унитарный, сохраняет спектр и энтропию \\
\quad\texttt{$s \leftarrow$ syndrome($\Gam$)} & // $3$-битная дуальная проверка $H$ (Прил.~\ref{app:fano}, Прин.~\ref{prin:selfdiag}) \\
\quad\texttt{$\Gam \leftarrow$ dissip\_fano($\Gam$, $\Delta\tau$)} & // $\gamma(\mathcal{P}_{\mathrm{Fano}}(\Gam)-\Gam)$ ($+$ опц.\ депол.) \\
\quad\texttt{$\Gam \leftarrow$ regen\_phiJ($\Gam$, $\Delta\tau$)} & // к $\varphi_J(\Gam)=(1-R)\mathcal{P}_{\mathrm{Fano}}(\Gam)+R\ket{\psi}\bra{\psi}$ \\
\quad\texttt{$\Gam \leftarrow$ U\_half($\Gam$, $\Delta\tau/2$)} & // симметрия Странга $+$ проекция PSD \\
\quad\texttt{assert invariants($\Gam$)} & // эрмитова, $\Tr{=}1$, PSD; точная $\Phi=(P-P_{\mathrm{diag}})/P_{\mathrm{diag}}$ \\
\quad\texttt{return ($\Gam$, $s$)} & \\
\bottomrule
\end{tabular}
\end{center}
Возвращённый $3$-битный синдром $s = He_j^T\pmod 2$ поступает на вход DIAKRISIS
на стороне \SYNARC{}; ненулевой $s$ однозначно локализует одиночный
деградировавший сектор (Приложение~\ref{app:fano}). Потиковая телеметрия
вычисляет точную $\Phi=(P-P_{\mathrm{diag}})/P_{\mathrm{diag}}$ (переходящую в
$7P-1$ \emph{только} на равномерной диагонали $\gamma_{kk}=1/7$), $R=1/(7P)$
(где $R\ge 1/3\Leftrightarrow P\le 3/7$ есть верхняя граница чистоты окна
жизнеспособности), $\CohE=(\gamma_{EE}^2+2\sum_{j\neq E}|\gamma_{Ej}|^2)/P$,
$D_{\mathrm{diff}}=1+6\,\CohE$ и $S_{\mathrm{vN}}=-\Tr(\Gam\log_2\Gam)\le\log_2 7$~бита.
Любое адаптивное демпфирование шага вблизи границы жизнеспособности есть
программный выбор управления (\status{D}), а не выведенная теорема.

\section{Глоссарий и перекрёстные ссылки на УГМ}
\label{app:glossary}

\paragraph{Термины.}
\begin{description}[leftmargin=0pt,itemsep=1pt]
  \item[\DM{}] регистр состояния: $\Gam\in\Dens(\mathbb{C}^7)$, $784$~байта ($48$ вещественных динамических параметров).
  \item[Тик] одно применение $\LindOmega$ (Ур.~\ref{eq:lomega}); единственная инструкция.
  \item[Линие-шина] одна из семи канонических прямых Зингера--Фано; единица интерконнекта (\S\ref{sec:fabric}).
  \item[Медиатор $k^\ast(i,j)$] единственный третий сектор прямой Фано для пары; детерминированная перемаршрутизация.
  \item[Коллапс] вывод интерконнекта, секторной двойственности Хэмминга, носителя инструкций и Фано-сжатой самомодели из единой инцидентности $\mathsf{N}$ (\S\ref{sec:collapse}).
  \item[Импеданс] накладные расходы эмуляции, которые субстрат добавляет для исполнения логической архитектуры (\S\ref{sec:realizability}).
  \item[Порт питания] неунитальный драйв $\Lind_{\mathrm{feed}}=r(\rho_{\mathrm{env}}-\Gam)$; шина импорта негэнтропии, разрыв которой ведёт к предсказуемому спаду к $\mathbb{1}/7$ (\S\ref{subsec:powerport}).
\end{description}

\paragraph{Использованные математические результаты и проектные параметры.}
Всякое логическое утверждение опирается на проверенную теорему (\status{Т}),
проектный параметр (\status{D}) или численную проверку в
\texttt{architecture/talos\_v0\_emulator.py},
\texttt{website/scripts/check\_core\_numbers.py} или
\texttt{website/research/operational-reconstruction/verify\_foundational\_bounds.py}.
\begin{center}\small
\begin{tabular}{@{}p{5.9cm}p{9.0cm}@{}}
\toprule
Результат / параметр & Использование в \TALOS{} \\
\midrule
Двоичное слияние плоскости и кода (\texttt{pg2\_5\_refutes\_N31\_hamming\_fusion}, \texttt{pg2EvenHammingUniqueSeven}) \status{Т} & единственность $N=7$ среди всех степеней простых $q\ge 2$ без тай-брейка (Теор.~\ref{thm:unique}), совпадение с $\Oct$ Гурвица \\
Дискретная ковариантность фабрики $\mathrm{PSL}(2,7)\ltimes U(1)^7$ \status{Т} & точная симметрия $\mathcal{P}_{\mathrm{Fano}}$ и $\mathcal{D}_{\mathrm{Fano}}$; непрерывная $G_2$ вращает базис (\S\ref{sec:collapse}, \S\ref{sec:platform}) \\
Правило отбора Фано / структурные константы $f_{ijk}$ \status{Т} & носитель умножения инструкции $=$ канонические прямые Зингера (L3) \\
Порождающая $M$ Хэмминга $[7,4,3]$ / дуальная проверочная $H$ симплекс-кода $[7,3,4]$ \status{Т} & секторная локализация отказов дистанции~$3$ ($HM^T\equiv 0\pmod 2$, \S\ref{sec:ecc}, Прил.~\ref{app:fano}) \\
Контракция Фано-сжатия $\Phi(\mathcal{P}_{\mathrm{Fano}}(\Gam))=\tfrac19\Phi(\Gam)$ \status{Т} & сохранение населённостей, внедиагонали $\times\tfrac13$; $\mathcal{D}_{\mathrm{Fano}}=\tfrac23\mathcal{D}_{\mathrm{atom}}$ (\S\ref{sec:fabric}) \\
Унитарная инвариантность спектра и энтропии \status{Т} & обратимый унитарный полушаг экспортирует $0$ логических бит; $S_{\mathrm{vN}}\le\log_2 7\approx 2{,}807$~бита (\S\ref{sec:energy}, Прил.~\ref{app:energy}) \\
Субстратная замкнутость (T-153), фазовый Gap (T-132) & классический комнатно-температурный $784$-Б регистр; лестница импеданса (\S\ref{sec:substrate}, \S\ref{sec:realizability}) \\
Жизнеспособность открытой системы (T-288/T-289) & унитальная динамика не растит $P$; неунитальный драйв/порт и аналитическая задержка отсечения (\S\ref{subsec:powerport}) \\
Каноническая заякоренная самомодель $\varphi_J(\Gam)$ & однопетлевая адаптация и тождество refresh--обучение (\S\ref{sec:selfmodel}, \S\ref{sec:platform}) \\
Вычислительная граница (BQP) & нет ускорения по классу сложности (честная граница, \S\ref{sec:intro}, \S\ref{sec:redteam}) \\
Точные наблюдаемые ($\Phi=(P-P_{\mathrm{diag}})/P_{\mathrm{diag}}$, $R=1/(7P)\ge 1/3\Leftrightarrow P\le 3/7$, $D_{\mathrm{diff}}=1+6\,\CohE$) \status{Т} & потиковая телеметрия жизнеспособности (\S\ref{sec:selfmodel}, \S\ref{sec:abi}, \S\ref{sec:platform}) \\
Проектная глубина композиции $d_{\max}=3$ (T-142) \status{D} & ограниченная $4$-регистровая башня ($\Cog_0$--$\Cog_3$, $3{,}136$~Б) $+$ горизонтальная экология (\S\ref{sec:composition}, \S\ref{sec:scaling}) \\
\bottomrule
\end{tabular}
\end{center}


\bibliographystyle{unsrtnat}
\bibliography{references}

\label{LastPage}
\end{document}
